\chapter{Differential equations: ODEs}
\label{vol02:ch07}

\epigraph{The laws of nature are but the mathematical thoughts of God.}{Attributed to Euclid by Leibniz; sourced here as an aphorism in the public domain rather than as a primary citation.}

\chapmeta{Half-life: foundational for the formal apparatus (the
canonical first- and second-order linear forms, the integrating
factor, the characteristic polynomial, the matrix exponential, the
phase plane); durable for the working numerical methods (Euler, RK4,
implicit Euler, the basic stiffness diagnosis); medium-life for the
specific solver tooling and library names of any given decade.
Archetypes: stability (the formal home; Lyapunov stability is
previewed here and developed in Volume X), control (the forced
response of the second-order linear oscillator previews Volumes
VIII-IX), transport (the heat-equation bridge to Chapter 12).
Exercise target: 45. Project track: simulation with analysis, per
Q55-vol02-ch07. Hazard class: none. Reader path default: standard,
with sections explicitly tagged. Named cases: Patriot missile system
at Dhahran (1991), invoked in the failure section as the prototype of
accumulated arithmetic error in long-running numerical integration;
Tacoma Narrows Bridge (1940), invoked in the forced-response section
to draw the line between resonance and aeroelastic flutter; Ariane 5
Flight 501 (1996), invoked in the failure section as the prototype of
range-versus-envelope failure in a numerical computation reused
outside the operating regime in which it was qualified.}

\noindent
A differential equation says how a quantity changes. The reader has
already met two specialised cases of the form: a derivative as a
rate, in Chapter 5, and an integral as a running accumulation, in
Chapter 6. An ordinary differential equation, an ODE, ties the two
together. It states that the rate of change of an unknown function
of one variable is a known function of the unknown and of the
variable. The chapter teaches the reader to recognise the canonical
forms (first-order linear, separable, second-order linear with
constant coefficients), to solve them by hand when a closed form
exists, to integrate them numerically when no closed form exists,
to read the qualitative behaviour off the phase plane, and to name
the failure modes by which an ODE solution either diverges, lies, or
silently shifts to the wrong attractor.

\section{What an ODE is, where it comes from}\pathtag{core}

An \engterm{ordinary differential equation} is an equation of the
form
\[
F\bigl(t,\, y(t),\, \dot{y}(t),\, \ddot{y}(t),\, \ldots,\, y^{(n)}(t)\bigr) = 0,
\]
where $y(t)$ is the unknown function, $t$ is the independent
variable, the dots and superscripts denote derivatives with respect
to $t$, and $F$ is a known function. The integer $n$ is the
\engterm{order} of the equation: the order of the highest derivative
that appears. The equation is \engterm{ordinary} because $y$ is a
function of one variable; if $y$ depended on more than one
independent variable, the equation would be a partial differential
equation, the subject of Chapter 12.

A \engterm{solution} of the ODE on an interval $I$ is a function
$y \colon I \to \R$ that is differentiable enough to make the
equation make sense, and that satisfies the equation pointwise on
$I$. An \engterm{initial-value problem (IVP)} is an ODE together
with $n$ side conditions $y(t_{0}) = y_{0}$, $\dot{y}(t_{0}) =
v_{0}$, and so on, that pin a single solution out of the family.
A \engterm{boundary-value problem (BVP)} replaces the $n$
conditions at one point with conditions distributed across two
points; BVPs arise routinely for second-order spatial problems and
return in the PDE chapter.

The conventional independent variable is $t$ for time and $x$ for
space, but the formal apparatus treats them identically. We use $t$
throughout this chapter for definiteness. The chapter on PDEs and
the chapter on multivariable calculus each treat space-dependent
problems in the same notation translated.

\subsection{Where ODEs come from}

ODEs appear whenever a balance law, a constitutive law, or a
controller is written down for a quantity that varies with one
parameter (typically time, sometimes position along an axis). The
canonical sources, all of which we revisit in this chapter or in
later volumes, are these.

\textbf{Newton's second law.} A particle of mass $m$ subject to a
force $F$ obeys $m \ddot{x} = F$. When $F$ depends on the position
$x$, the velocity $\dot{x}$, and time $t$, the equation is a
second-order ODE in $x(t)$. A spring force $-k x$, a viscous force
$-c \dot{x}$, and an external drive $F(t)$ together produce
$m \ddot{x} + c \dot{x} + k x = F(t)$, the canonical form whose
analysis occupies §7.4 and §7.5.

\textbf{The RC circuit.} A resistor $R$ in series with a capacitor
$C$ across a voltage $V(t)$ obeys, by Kirchhoff's voltage law and
the constitutive laws $V_{R} = R i$ and $V_{C} = q/C$ with $i =
\dot{q}$, the equation $R \dot{q} + q/C = V(t)$. The same equation
in terms of capacitor voltage $v_{C} = q/C$ reads $RC \dot{v}_{C} +
v_{C} = V(t)$. The product $\tau = RC$ has dimensions of time and
sets the circuit's response time. The discharge problem
$RC \dot{v}_{C} + v_{C} = 0$ is the prototype first-order linear
ODE; its solution motivates the entire Chapter.

\textbf{Population growth.} A population whose per-capita
reproduction rate is constant obeys $\dot{N} = r N$, the prototype
of exponential growth. With finite resources, the rate decreases
with population size: $\dot{N} = r N (1 - N/K)$, the logistic
equation. Both are first-order ODEs; the first is linear in $N$, the
second is not.

\textbf{Mass-balance.} A tank receives a substance at known
concentration and rate, mixes it perfectly, and discharges at a
volume rate that may match or differ from the inflow. Conservation
of mass for the dissolved species gives $\dot{m} = (\text{inflow
rate}) - (\text{outflow rate})$, where each term is a function of
$m(t)$ and $t$. Mass-balance is the most common applied source of
first-order ODEs across chemical, environmental, biomedical, and
process engineering. Drug pharmacokinetics, lake pollution
clearance, well-stirred reactor models, and household water-heater
dynamics are all mass-balance ODEs in disguise.

\begin{archetype}[Stability, formally introduced here]
The stability archetype was previewed in Chapter 3 (rotation as
structure-preserving) and named in Volume I as the recurring
question: when a system is perturbed, does it return to where it
was, drift away, or settle into a new state? In ODE language the
question is asked of \engterm{equilibria}: points $y^{*}$ where
$\dot{y} = f(y^{*}) = 0$. An equilibrium is \engterm{stable} (in
the sense of Lyapunov) if a small perturbation produces a small
deviation, and \engterm{asymptotically stable} if the perturbation
decays back to zero. The chapter develops the criterion through
linearisation in §7.7. The full Lyapunov-function machinery, by
which stability is certified without solving the equation, is
previewed in §7.7 and developed at working depth in Volume X.
\end{archetype}

\subsection{Existence, uniqueness, and what an engineer needs from them}\pathtag{standard}

The Picard-Lindel\"of theorem says that the IVP $\dot{y} = f(t, y)$
with $y(t_{0}) = y_{0}$ has a unique solution on an interval around
$t_{0}$ provided $f$ is continuous in $t$ and Lipschitz continuous
in $y$ near $(t_{0}, y_{0})$. The proof, by the contraction-mapping
theorem applied to the integral equation
$y(t) = y_{0} + \int_{t_{0}}^{t} f(s, y(s))\,ds$, lives in a course
on real analysis. The engineering reader does not need the proof.
The reader needs three working consequences.

First, the standard right-hand sides of physics, chemistry, and
biology (smooth functions of $y$ and $t$) satisfy the hypotheses on
their working domain, and the IVP has a unique solution.

Second, when the right-hand side fails to be Lipschitz, uniqueness
can fail. The textbook example is $\dot{y} = \sqrt{y}$ with $y(0) =
0$: the trivial solution $y \equiv 0$ and the family of
late-launching solutions $y(t) = (t - t_{0})^{2}/4$ for any $t_{0}
\geq 0$ all satisfy the equation and the initial condition. A
real-world echo: at the moment a lubricated bearing transitions
between regimes, the friction law has a kink and an idealised model
loses uniqueness. Engineers respond by smoothing the model,
restricting the operating regime, or treating the transition as a
switching event.

Third, even when uniqueness holds locally, solutions can fail to
extend to all $t$. The equation $\dot{y} = y^{2}$ with $y(0) = 1$
has the solution $y(t) = 1/(1 - t)$, which blows up at $t = 1$.
Models that produce finite-time blow-up typically signal that the
modelled mechanism (combustion, contagion in an unbounded domain,
gravitational collapse) saturates physically before the model says
so. The model is correct on its working interval and wrong about
the singularity; the engineer's job is to know the working
interval.

\begin{estimation}
\textbf{Question.} An RC circuit with $R = 10\,\si{\kilo\ohm}$ and
$C = 100\,\si{\micro\farad}$ is charged to $5\,\si{\volt}$ and
allowed to discharge through the resistor. Estimate the time to
fall to half its initial voltage and the time to fall below
$50\,\si{\milli\volt}$, both before reaching for the formal
solution.

\smallskip
\textit{Estimate, before reading on.}

\smallskip
The time scale of an RC circuit is $\tau = RC = 10^{4}\,\si{\ohm}
\times 10^{-4}\,\si{\farad} = 1\,\si{\second}$. An exponential
discharge falls by a factor of $e \approx 2.72$ per time constant.
To fall by a factor of $2$ takes $\ln 2 \approx 0.69$ time
constants, so the half-time is about $0.69\,\si{\second}$. To fall
from $5\,\si{\volt}$ to $50\,\si{\milli\volt}$ is a factor-of-100
drop, or $\ln 100 \approx 4.6$ time constants, so about
$4.6\,\si{\second}$. The estimate sets expectations: the
oscilloscope trace settles within five seconds, and a measurement
that wants four-digit precision on the residual voltage should
sample for at least seven time constants ($e^{-7} \approx 10^{-3}$).

The formal solution in §7.2 will give $v(t) = 5 e^{-t/\tau}$ and
confirm both estimates. The estimation is what an engineer does
before fitting an oscilloscope trace; the calculation is what
follows.
\end{estimation}

\subsection{Linearity and the working classification}

The chapter's working taxonomy is small. An ODE is \engterm{linear}
if it can be written as a linear combination of $y$ and its
derivatives, with coefficients that depend on $t$ but not on $y$:
\[
a_{n}(t)\, y^{(n)} + a_{n-1}(t)\, y^{(n-1)} + \cdots + a_{0}(t)\, y = q(t).
\]
The function $q(t)$ is the \engterm{forcing}; when $q \equiv 0$ the
equation is \engterm{homogeneous}. When the coefficients $a_{i}$
do not depend on $t$, the equation has \engterm{constant
coefficients}; otherwise it has variable coefficients. An ODE that
is not linear is \engterm{nonlinear}. A first-order ODE that can be
written $\dot{y} = g(t)\, h(y)$ is \engterm{separable}; separable
equations may be linear or nonlinear.

The classification matters because the solution methods follow it.
Linear equations have a structured solution theory (superposition of
homogeneous solutions plus a particular solution); separable
equations integrate by quadrature; nonlinear non-separable equations
generally require numerical methods.

\section{First-order linear ODEs}\pathtag{core}

The standard form of a first-order linear ODE is
\[
\dot{y} + p(t)\, y = q(t),
\]
where $p$ and $q$ are known continuous functions of $t$ on the
interval of interest. Every first-order linear ODE can be brought
to this form by dividing through by the coefficient of $\dot{y}$.

\subsection{The integrating factor}

The trick that solves the standard form is an old one: multiply
both sides by an as-yet-unknown function $\mu(t)$, and choose $\mu$
so that the left-hand side becomes the derivative of a product.
Recall the product rule: $\frac{d}{dt}(\mu y) = \dot{\mu} y + \mu
\dot{y}$. Multiplying the equation by $\mu$ gives
\[
\mu \dot{y} + \mu p\, y = \mu q.
\]
For the left-hand side to equal $\frac{d}{dt}(\mu y) = \dot{\mu} y +
\mu \dot{y}$, we need $\dot{\mu} = \mu p$, an ODE for $\mu$ alone
that we can solve by inspection: $\mu(t) = \exp\bigl(\int p(t)\,dt\bigr)$.
We choose any antiderivative; the constant of integration cancels
in the next step. The function
\[
\mu(t) = \exp\!\left(\int p(t)\,dt\right)
\]
is the \engterm{integrating factor}.

With $\mu$ in hand, the equation becomes
\[
\frac{d}{dt}(\mu(t)\, y(t)) = \mu(t)\, q(t),
\]
which integrates directly:
\[
\mu(t)\, y(t) = \int \mu(t)\, q(t)\,dt + C,
\]
or, written for the IVP with $y(t_{0}) = y_{0}$,
\[
y(t) = \frac{1}{\mu(t)}\!\left[\mu(t_{0})\, y_{0} + \int_{t_{0}}^{t} \mu(s)\, q(s)\,ds\right].
\]
This is the closed-form solution to every first-order linear IVP,
provided the integral on the right can be evaluated. Whether it can
or cannot is a question about $p$ and $q$, not about the method.

\subsection{The RC discharge worked through}

Take the discharge problem $RC \dot{v} + v = 0$ with $v(0) =
V_{0}$. Bring to standard form by dividing through by $RC$:
$\dot{v} + (1/RC) v = 0$. Here $p(t) = 1/(RC) = 1/\tau$ is
constant; $q(t) = 0$.

The integrating factor is $\mu(t) = \exp(t/\tau)$. The right-hand
integral is zero. The IVP solution reduces to $\mu(t) v(t) =
\mu(0) V_{0}$, or
\[
v(t) = V_{0}\, e^{-t/\tau}.
\]
Substituting $V_{0} = 5\,\si{\volt}$ and $\tau = 1\,\si{\second}$
recovers $v(0.69) = 2.5\,\si{\volt}$ and $v(4.6) =
0.050\,\si{\volt}$, matching the estimation.

\subsection{The driven RC circuit}

Now drive the circuit with a constant input $V_{0}$ starting at
$t = 0$, with $v(0) = 0$ (capacitor initially empty). The equation
is $RC \dot{v} + v = V_{0}$ for $t > 0$. Standard form:
$\dot{v} + (1/\tau) v = V_{0}/\tau$. Integrating factor $\mu =
e^{t/\tau}$. Multiplying:
\[
\frac{d}{dt}(e^{t/\tau} v) = \frac{V_{0}}{\tau}\, e^{t/\tau}.
\]
Integrating from $0$ to $t$:
\[
e^{t/\tau} v(t) - v(0) = V_{0}\bigl(e^{t/\tau} - 1\bigr),
\]
so
\[
v(t) = V_{0}\bigl(1 - e^{-t/\tau}\bigr).
\]
The solution rises from $0$ toward $V_{0}$, reaching $63\,\%$ of
$V_{0}$ at one time constant, $86\,\%$ at two, $95\,\%$ at three,
$99\,\%$ at five. The reader who works with first-order systems
should know these percentages cold; they are the working calibration
of "how long is long enough" in any process whose response is set
by a single time constant.

% Figure: RC circuit transient. The charging and discharging curves
% on a normalised axis (time in units of tau, voltage as fraction of V0).
\begin{figure}[htbp]
\centering
\begin{tikzpicture}[
  axis/.style={draw=engblack, -{Stealth}, thick},
  charge/.style={draw=engblue, very thick},
  discharge/.style={draw=engred, very thick},
  guide/.style={draw=enggray, dashed, thin},
  label/.style={font=\small},
  domain=0:5, samples=200,
]
% Axes
\draw[axis] (0,0) -- (10.5,0) node[anchor=west, label] {$t/\tau$};
\draw[axis] (0,0) -- (0,4.8)  node[anchor=south, label] {$v/V_{0}$};

% Tick marks on x-axis
\foreach \x/\m in {0/0, 2/1, 4/2, 6/3, 8/4, 10/5} {
  \draw (\x,-0.08) -- (\x,0.08);
  \node[font=\scriptsize, below] at (\x,-0.08) {\m};
}
% Tick marks on y-axis
\foreach \y/\h in {0/0, 1/0.25, 2/0.50, 3/0.75, 4/1.00} {
  \draw (-0.08,\y) -- (0.08,\y);
  \node[font=\scriptsize, left] at (-0.08,\y) {\h};
}

% Charging curve: y = 4 * (1 - exp(-x/2))   (4 units == 1.00)
\draw[charge] plot[domain=0:10, samples=200]
  ({\x}, {4*(1 - exp(-\x/2))});
% Discharging curve: y = 4 * exp(-x/2)
\draw[discharge] plot[domain=0:10, samples=200]
  ({\x}, {4*exp(-\x/2)});

% Reference horizontal lines at 0.63 and 0.37
\draw[guide] (0, 2.52) -- (10, 2.52);
\node[engblack, font=\scriptsize, anchor=west] at (10.05, 2.52) {$0.63$};
\draw[guide] (0, 1.48) -- (10, 1.48);
\node[engblack, font=\scriptsize, anchor=west] at (10.05, 1.48) {$0.37$};

% Vertical guide at one time constant (x=2)
\draw[guide] (2, 0) -- (2, 4);
\node[font=\scriptsize, below right, enggray] at (2, 0.5) {$t=\tau$};

% Labels on curves
\node[engblue, font=\small, anchor=south west] at (5.5, 3.6)
  {charging $v=V_{0}(1-e^{-t/\tau})$};
\node[engred, font=\small, anchor=north west] at (5.5, 0.4)
  {discharging $v=V_{0}e^{-t/\tau}$};

\end{tikzpicture}
\caption[RC circuit transient on the normalised time axis]{First-order
RC transients on the normalised axis $t/\tau$. The charging curve
(blue, from $0$ to $V_{0}$) and the discharging curve (red, from
$V_{0}$ to $0$) cross the $0.63 V_{0}$ and $0.37 V_{0}$ lines at one
time constant. The engineer who has internalised the percentage
calibration ($63\,\%$ at $1\tau$, $86\,\%$ at $2\tau$, $95\,\%$ at
$3\tau$, $99\,\%$ at $5\tau$) reads first-order settling time off
this picture by eye. Computed from $v=V_{0}(1-e^{-t/\tau})$ and
$v=V_{0}e^{-t/\tau}$.}
\label{fig:vol02:ch07:rc-transient}
\end{figure}


\subsection{Linearity in action: superposition}

The general first-order linear equation $\dot{y} + p y = q$ has
solution $y = y_{h} + y_{p}$, where $y_{h}$ is any solution of the
homogeneous equation $\dot{y} + p y = 0$ and $y_{p}$ is any
particular solution of the full equation. The homogeneous solutions
form a one-parameter family (multiples of $1/\mu(t)$); the
particular solution carries the response to $q$. The
initial condition pins the constant in front of $y_{h}$. This
structure is the entire content of "linear" for first-order
equations and recurs at higher order.

\subsection{Variable coefficients: a worked example}

Solve $\dot{y} + 2 y / t = t$ on $t > 0$ with $y(1) = 0$.

In standard form, $p(t) = 2/t$ and $q(t) = t$. The integrating
factor is
\[
\mu(t) = \exp\!\left(\int \frac{2}{t}\,dt\right) = \exp(2 \ln t) = t^{2}.
\]
The equation becomes $\frac{d}{dt}(t^{2} y) = t^{3}$, integrating
to $t^{2} y = t^{4}/4 + C$. The initial condition $y(1) = 0$ gives
$C = -1/4$. The solution is
\[
y(t) = \frac{t^{2}}{4} - \frac{1}{4 t^{2}}.
\]
A check: at $t = 1$, $y = 0$. The derivative is $\dot{y} = t/2 +
1/(2 t^{3})$. Substituting back, $\dot{y} + 2 y/t = (t/2 + 1/(2
t^{3})) + 2(t^{2}/4 - 1/(4 t^{2}))/t = t/2 + 1/(2 t^{3}) + t/2 -
1/(2 t^{3}) = t$. The solution checks.

The reader should always check a solved ODE by substitution. It
costs ten seconds and catches sign errors that would otherwise
travel into the next calculation.

\section{Separable equations}\pathtag{core}

A first-order ODE is \engterm{separable} if it can be written
$\dot{y} = g(t) h(y)$. Separation rearranges to
$\frac{dy}{h(y)} = g(t)\,dt$ and integrates both sides:
\[
\int \frac{dy}{h(y)} = \int g(t)\,dt + C.
\]
The result is an implicit equation for $y$ as a function of $t$;
solving for $y$ explicitly is sometimes possible and sometimes not.

\subsection{Logistic growth worked through}

The logistic equation $\dot{N} = r N (1 - N/K)$ is separable, with
$g(t) = r$ (constant) and $h(N) = N (1 - N/K)$. Separate and
integrate:
\[
\int \frac{dN}{N (1 - N/K)} = \int r\,dt.
\]
The left-hand integrand decomposes by partial fractions:
$\frac{1}{N (1 - N/K)} = \frac{1}{N} + \frac{1/K}{1 - N/K}$.
Integrating gives $\ln |N| - \ln |1 - N/K| = r t + C$, or, after
exponentiation and absorbing the constant,
\[
\frac{N}{1 - N/K} = A\, e^{r t},
\]
with $A$ determined by $N(0) = N_{0}$. Solving for $N$:
\[
N(t) = \frac{K\, N_{0}\, e^{r t}}{K + N_{0}(e^{r t} - 1)}.
\]
The solution rises from $N_{0}$ toward $K$ as $t \to \infty$. The
inflection point, where the growth rate is maximised, is at $N =
K/2$. The reader who fits a logistic model to data should know
this property: the data point at peak rate identifies $K/2$
directly.

% Figure: slope field for the logistic equation, with two integral
% curves shown to demonstrate convergence to the carrying capacity.
\begin{figure}[htbp]
\centering
\begin{tikzpicture}[
  axis/.style={draw=engblack, -{Stealth}, thick},
  field/.style={draw=enggray, thin, -{Latex[length=1.4mm, width=1.0mm]}},
  curveA/.style={draw=engblue, very thick},
  curveB/.style={draw=engred, very thick},
  curveC/.style={draw=engamber, very thick},
  label/.style={font=\small},
]
% Axes (scaled: x from 0 to 10 (time), y from 0 to 6 (N/K, where K=1 -> 4 units))
\draw[axis] (0,0) -- (11,0) node[anchor=west, label] {$t$};
\draw[axis] (0,0) -- (0,5.4) node[anchor=south, label] {$N$};

% Ticks
\foreach \x in {0, 2, 4, 6, 8, 10} {
  \draw (\x,-0.07) -- (\x,0.07);
  \node[font=\scriptsize, below] at (\x,-0.07) {\x};
}
\foreach \y/\h in {0/0, 1/0.25, 2/0.50, 3/0.75, 4/1.00} {
  \draw (-0.07,\y) -- (0.07,\y);
  \node[font=\scriptsize, left] at (-0.07,\y) {\h};
}

% Carrying capacity line K=1 (at y=4)
\draw[draw=enggray, dashed, thin] (0,4) -- (11,4);
\node[font=\scriptsize, enggray, anchor=west] at (10.5, 4.2) {$K$};

% Slope field: short arrows; dN/dt = N(1 - N/4) (r=1, K=4 in plot units)
% Plot units: y in 0..4 is N in 0..1
\foreach \tx in {0.4, 1.2, 2.0, 2.8, 3.6, 4.4, 5.2, 6.0, 6.8, 7.6, 8.4, 9.2, 10.0} {
  \foreach \ny in {0.2, 0.6, 1.0, 1.4, 1.8, 2.2, 2.6, 3.0, 3.4, 3.8, 4.2, 4.6, 5.0} {
    % slope = N(1 - N/4) where N = ny/4 maps to actual N
    % Use ny as the rendered y; slope_render = (ny/4)*(1 - ny/4)
    \pgfmathsetmacro{\slp}{(\ny/4)*(1 - \ny/4)}
    \pgfmathsetmacro{\norm}{sqrt(1 + (\slp)^2)}
    \pgfmathsetmacro{\dx}{0.35/\norm}
    \pgfmathsetmacro{\dy}{0.35*\slp/\norm}
    \draw[field] (\tx - \dx, \ny - \dy) -- (\tx + \dx, \ny + \dy);
  }
}

% Three integral curves: logistic solutions with K=1, r=1, three N0
% Plot uses N_plot = 4 * K * N0 * exp(t) / (K + N0 (exp(t) - 1))
% A: N0 = 0.05; B: N0 = 0.5; C: N0 = 1.5 (overshoot)
\draw[curveA] plot[domain=0:6, samples=80]
  ({\x}, {4 / (1 + ((1 - 0.05)/0.05) * exp(-\x))});
\draw[curveB] plot[domain=0:6, samples=80]
  ({\x}, {4 / (1 + ((1 - 0.5)/0.5) * exp(-\x))});
\draw[curveC] plot[domain=0:6, samples=80]
  ({\x}, {4 / (1 + ((1 - 1.5)/1.5) * exp(-\x))});

% Labels
\node[engblue, font=\scriptsize, anchor=west] at (4.5, 3.0) {$N_{0}=0.05\,K$};
\node[engred, font=\scriptsize, anchor=west] at (3.2, 4.45) {$N_{0}=0.5\,K$};
\node[engamber, font=\scriptsize, anchor=west] at (1.0, 5.1) {$N_{0}=1.5\,K$ (over)};

\end{tikzpicture}
\caption[Slope field for the logistic equation]{Slope field for the
logistic equation $\dot{N}=rN(1-N/K)$ with $r=1$ and $K=1$, with
three integral curves drawn at initial conditions $N_{0}=0.05\,K$,
$N_{0}=0.5\,K$, and $N_{0}=1.5\,K$. All solutions converge to the
carrying capacity $K$. The inflection point at $N=K/2$ marks the
maximum growth rate; trajectories below approach the inflection
from below and accelerate, trajectories above decay monotonically.
The slope field summarises every solution at once; the integral
curves are particular solutions selected by initial condition.}
\label{fig:vol02:ch07:slope-field}
\end{figure}


\subsection{Separable but not closed-form}

The equation $\dot{y} = e^{-t^{2}}/(1 + y^{2})$ is separable, with
$h(y) = 1/(1 + y^{2})$ and $g(t) = e^{-t^{2}}$. Separation gives
$(1 + y^{2})\,dy = e^{-t^{2}}\,dt$. The left integrates to $y +
y^{3}/3$; the right does not have an elementary antiderivative
(the error function $\operatorname{erf}$ is not elementary, and
$\int e^{-t^{2}}\,dt = (\sqrt{\pi}/2) \operatorname{erf}(t)$ is the
canonical example of an integral the reader can name but not
elementary-evaluate). The solution exists in closed form once
$\operatorname{erf}$ is treated as known; it does not exist in
elementary closed form.

The lesson recurs throughout the chapter: separability is a
structural property that promises a quadrature; whether the
quadrature evaluates to elementary functions is a separate
question.

\section{Second-order linear with constant coefficients}\pathtag{core}

The canonical second-order linear ODE with constant coefficients is
\[
m \ddot{x} + c \dot{x} + k x = F(t),
\]
where $m, c, k$ are positive constants. Physically, $m$ is mass, $c$
is damping (in $\si{\newton\second\per\meter}$), $k$ is stiffness
(in $\si{\newton\per\meter}$), and $F$ is the external forcing
(force in $\si{\newton}$). The form covers a mass-spring-dashpot
system, a series RLC circuit (with $L$ playing the role of $m$, $R$
of $c$, $1/C$ of $k$), the small-amplitude pendulum, and any
linearised second-order system the engineer encounters.

\subsection{The homogeneous problem and the characteristic polynomial}

Set $F = 0$ first. The homogeneous equation
$m \ddot{x} + c \dot{x} + k x = 0$ has solutions of the form $x(t) =
e^{\lambda t}$ for some constant $\lambda$, on the strength of the
fact that $\frac{d}{dt}e^{\lambda t} = \lambda e^{\lambda t}$. Plug
in:
\[
m \lambda^{2} e^{\lambda t} + c \lambda e^{\lambda t} + k e^{\lambda t} = e^{\lambda t}(m \lambda^{2} + c \lambda + k) = 0.
\]
Since $e^{\lambda t}$ is never zero, the equation reduces to the
\engterm{characteristic polynomial}
\[
m \lambda^{2} + c \lambda + k = 0,
\]
with roots
\[
\lambda_{\pm} = \frac{-c \pm \sqrt{c^{2} - 4 m k}}{2 m}.
\]
The discriminant $c^{2} - 4 m k$ classifies the homogeneous
behaviour into three regimes.

\textbf{Overdamped} ($c^{2} > 4 m k$): two distinct real negative
roots; the solution is $x(t) = A_{1} e^{\lambda_{1} t} + A_{2}
e^{\lambda_{2} t}$, a sum of two decaying exponentials. The system
returns to equilibrium without oscillating. Worst case for response
time; useful when no overshoot is acceptable.

\textbf{Critically damped} ($c^{2} = 4 m k$): a repeated real root
$\lambda = -c/(2 m)$; the solution is $x(t) = (A_{1} + A_{2} t)\,
e^{\lambda t}$. Fastest non-oscillating return to equilibrium.
Standard target in instrument and door-closer design.

\textbf{Underdamped} ($c^{2} < 4 m k$): two complex conjugate roots
$\lambda = -\zeta \omega_{0} \pm i \omega_{d}$, where
$\omega_{0} = \sqrt{k/m}$ is the \engterm{natural frequency}, $\zeta
= c / (2 \sqrt{m k})$ is the \engterm{damping ratio}, and
$\omega_{d} = \omega_{0}\sqrt{1 - \zeta^{2}}$ is the
\engterm{damped frequency}. The solution is
\[
x(t) = e^{-\zeta \omega_{0} t} \bigl(A \cos(\omega_{d} t) + B \sin(\omega_{d} t)\bigr),
\]
an exponentially decaying sinusoid. Most physical oscillators
(pendulum, vibrating string, AC circuit) sit in this regime.

The form $\omega_{0}, \zeta$ is the engineer's working notation.
The reader should be able to read off $\omega_{0}$ and $\zeta$ from
$m, c, k$ by inspection.

\subsection{Worked example: the underdamped oscillator}

Take $m = 1\,\si{\kilogram}$, $c = 0.5\,\si{\newton\second\per\meter}$,
$k = 100\,\si{\newton\per\meter}$. Then $\omega_{0} = \sqrt{k/m} =
10\,\si{\radian\per\second}$, and $\zeta = c / (2\sqrt{m k}) = 0.5
/ 20 = 0.025$. The damping ratio is small; the system is
lightly damped. The damped frequency is $\omega_{d} = 10
\sqrt{1 - 0.025^{2}} \approx 9.997\,\si{\radian\per\second}$, very
close to $\omega_{0}$ for light damping. The decay rate is $\zeta
\omega_{0} = 0.25\,\si{\per\second}$, giving a $1/e$ amplitude
decay time of $4\,\si{\second}$, or about $40$ oscillation periods
($T = 2\pi/\omega_{d} \approx 0.628\,\si{\second}$).

A bell, a tuning fork, and a high-quality pendulum clock all sit in
this lightly damped regime. The \engterm{quality factor} $Q = 1 /
(2 \zeta) = 20$ counts the oscillation periods to amplitude decay
of $e^{-\pi}$; $Q = 20$ is modest. A laboratory tuning fork has $Q
\sim 10^{4}$.

\subsection{The complex-exponential representation}

The complex-exponential form $x(t) = \operatorname{Re}\bigl(\hat{x}\,
e^{\lambda t}\bigr)$ with $\hat{x}$ complex is more compact than
the cosine-plus-sine form and is the working notation for the rest
of the chapter and Volumes III, VIII, and IX. The reader who has
worked Chapter 3 on complex phasors should recognise the apparatus
without strain. The two real parameters $A, B$ become one complex
parameter $\hat{x}$; the cosine and sine become the real and
imaginary parts of $e^{i \omega_{d} t}$.

\section{Forced response: resonance, beating}\pathtag{core}

Set $F(t) = F_{0} \cos(\omega t)$ in $m \ddot{x} + c \dot{x} + k x =
F(t)$. The full solution is $x = x_{h} + x_{p}$, with $x_{h}$ the
homogeneous solution (decaying exponential or decaying sinusoid)
and $x_{p}$ the steady-state response to the drive.

Try $x_{p}(t) = A \cos(\omega t - \varphi)$ with amplitude $A$ and
phase lag $\varphi$ to be found. Substitute and use the
trigonometric expansions; or, more cleanly, work in complex
exponentials. With $F(t) = \operatorname{Re}(F_{0}\, e^{i \omega
t})$ and $x_{p}(t) = \operatorname{Re}(\hat{X}\, e^{i \omega t})$:
\[
m (i\omega)^{2} \hat{X} + c (i\omega) \hat{X} + k \hat{X} = F_{0},
\]
giving
\[
\hat{X} = \frac{F_{0}}{(k - m \omega^{2}) + i c \omega}.
\]
The amplitude is the modulus, the phase lag is the argument:
\[
A(\omega) = \frac{F_{0}}{\sqrt{(k - m\omega^{2})^{2} + (c\omega)^{2}}},
\quad
\tan\varphi(\omega) = \frac{c\omega}{k - m\omega^{2}}.
\]

\subsection{Resonance}

The amplitude $A(\omega)$ has a maximum at the \engterm{resonant
frequency} $\omega_{r} = \omega_{0}\sqrt{1 - 2\zeta^{2}}$, slightly
below $\omega_{0}$ for nonzero damping. At resonance, the amplitude
is approximately
\[
A(\omega_{r}) \approx \frac{F_{0}}{c \omega_{0}} = \frac{F_{0}}{k}\, \frac{1}{2\zeta\sqrt{1-\zeta^{2}}} \approx \frac{F_{0}}{k} Q
\]
for light damping ($\zeta \ll 1$), where $Q$ is the quality factor.
The static-deflection amplitude $F_{0}/k$ is multiplied by $Q$ at
resonance: a $Q = 100$ system has $100\times$ static deflection at
its resonance.

The phase lag is $\pi/2$ at $\omega = \omega_{0}$. At low frequency
the response follows the drive in phase; at high frequency the
response lags the drive by $\pi$ (the inertia dominates); at the
natural frequency the phase lag is exactly a quarter cycle.

The Tacoma Narrows Bridge collapse of 7 November 1940 is the
canonical popular example of resonance, although the actual
mechanism was aeroelastic flutter, a more complicated
fluid-structure coupling that is not steady-state forced response
of a linear oscillator. The Federal Works Agency report, prepared
by Ammann, von K\'arm\'an, and Woodruff and released on 28 March
1941, identified the cause as torsional self-excited oscillation
of the deck section at a steady wind of approximately
$19\,\si{\meter\per\second}$: at the critical wind speed, the
angular displacement of the deck modified the aerodynamic moment in
a manner that fed energy back into the torsional motion
\cite{acc:tacoma-narrows-1941}. The deck had been designed for
static wind loads with no aerodynamic analysis of the cross-section
as a body in steady flow. The registry entry
\repopath{docs/research/accidents/tacoma-narrows-1940.md} carries
the careful discussion. Engineers do not invoke "resonance" as a
shorthand for "vibration that grew." Resonance is a specific
amplitude peak in the frequency response of a linear oscillator;
flutter, parametric instability, and limit cycles are distinct
mechanisms with distinct mathematics. The distinction has practical
weight: a resonance is mitigated by adding damping or detuning the
forcing; a flutter is mitigated by changing the cross-section
geometry or the structural stiffness. Flutter, parametric
instability, and limit cycles are taken up in Volumes III, VIII,
and X.

\subsection{Beating}

When two sinusoids of nearby frequencies $\omega_{1}$ and
$\omega_{2}$ add, the result is a fast oscillation at the average
frequency $(\omega_{1} + \omega_{2})/2$ modulated by a slow envelope
at the difference frequency $|\omega_{1} - \omega_{2}|/2$. The
identity
\[
\cos(\omega_{1} t) + \cos(\omega_{2} t) = 2 \cos\!\left(\frac{\omega_{1} + \omega_{2}}{2} t\right) \cos\!\left(\frac{\omega_{1} - \omega_{2}}{2} t\right)
\]
is the trigonometric content. \engterm{Beating} is the audible
signature: a piano tuner adjusting two strings to the same
fundamental hears the beat slow and stop as the strings come into
tune. In a forced oscillator with drive frequency near the natural
frequency, the homogeneous transient and the steady-state response
beat against each other for several time constants of damping
before the transient decays.

\begin{estimation}
\textbf{Question.} A wine glass has a fundamental rim mode near
$f_{0} = 700\,\si{\hertz}$ and rings audibly for about
$3\,\si{\second}$ after a flick. Estimate the steady-state
amplitude of rim motion when a singer drives the glass at exactly
its natural frequency with a sound-pressure level of
$120\,\si{\decibel}$ (roughly the loudest a trained voice
sustains at close range). Estimate the threshold pressure at
which the rim strain exceeds the glass's working strain limit
($\approx 10^{-3}$). Estimate before computing.

\smallskip
\textit{Estimate, before reading on.}

\smallskip
The audible ring of about $3\,\si{\second}$ is roughly a $1/e$
amplitude time. The angular natural frequency is $\omega_{0} = 2
\pi f_{0} \approx 4400\,\si{\radian\per\second}$. The damping ratio
follows from $\zeta \omega_{0} \tau_{\text{ring}} \approx 1$, so
$\zeta \approx 1/(\omega_{0} \tau) \approx 1/(4400 \times 3)
\approx 7.5 \times 10^{-5}$, giving a quality factor $Q = 1/(2
\zeta) \approx 6500$. A sound-pressure level of $120\,\si{\decibel}$
corresponds to a pressure amplitude of about $20\,\si{\pascal}$
(the reference $20\,\si{\micro\pascal}$ scaled by $10^{120/20}$).
Over a glass rim of effective area $\sim 5\,\si{\centi\meter
\squared} = 5 \times 10^{-4}\,\si{\meter\squared}$, the driving
force is $F_{0} \approx 20 \times 5 \times 10^{-4} = 10^{-2}\,
\si{\newton}$. The static rim deflection $F_{0}/k$ requires an
order-of-magnitude estimate of the rim stiffness; for a thin-walled
goblet, $k \sim 10^{4}\,\si{\newton\per\meter}$ is in the right
range, so $F_{0}/k \sim 10^{-6}\,\si{\meter} = 1\,\si{\micro
\meter}$. At resonance the response is amplified by $Q$, giving a
peak rim displacement of $Q \times 10^{-6}\,\si{\meter} \approx 7
\,\si{\milli\meter}$. For a rim of radius $\sim 4\,\si{\centi
\meter}$, this displacement corresponds to a strain on the order
of $7 \times 10^{-3}/4 \times 10^{-2} = 0.18 \approx 18\,\%$,
two orders of magnitude above the working strain limit of glass.
The estimate predicts the glass shatters. The phenomenon is
real; trained singers have broken wine glasses on stage, and the
mechanism is exactly the $Q$-fold amplification of a low static
deflection by a high-$Q$ resonance.

The estimate also explains why a glass cannot be broken at off-
resonance frequencies: a $20\,\si{\hertz}$ detuning from
$700\,\si{\hertz}$ takes the drive well outside the half-power
bandwidth $\omega_{0}/Q \approx 0.7\,\si{\radian\per\second}$
($\approx 0.1\,\si{\hertz}$), and the amplification collapses to
the static deflection. Engineering practice for resonant
sensors and actuators lives entirely inside this trade-off: a
narrow bandwidth buys amplification but demands frequency
control to within the bandwidth's own width.
\end{estimation}

\subsection{Quality factor and bandwidth}

The frequency response $A(\omega)$ has a half-power bandwidth
$\Delta\omega = \omega_{0}/Q$ (approximately, for light damping).
The quality factor is the working currency of resonant-system
design: a high-$Q$ filter is selective in frequency but slow to
settle (settling time $\sim Q/\omega_{0}$); a low-$Q$ filter is
broad but fast. The trade-off is set by the physics, not by the
engineer's preference. A radio receiver, a laser cavity, an
electrochemical impedance spectrometer, a structural damper, and a
mechanical vibration isolator all live inside this trade-off.

\begin{principle}[The forced linear oscillator is the prototype of control]
A second-order linear oscillator under sinusoidal drive admits an
exact closed-form steady-state response with three working
parameters: the natural frequency $\omega_{0}$, the damping ratio
$\zeta$ (or the quality factor $Q = 1/(2\zeta)$), and the
forcing-to-stiffness ratio $F_{0}/k$. Every control system that
asks "how does the output follow a reference signal" begins with
this prototype. The frequency response, the phase lag, the
resonance peak, the bandwidth, and the trade-off between selectivity
and settling time are the working vocabulary of Volumes VIII (Signals
and control) and IX (Cyber-physical systems). The reader who has
internalised the prototype will read those volumes as elaborations
of a model already in hand.
\end{principle}

\section{Systems of ODEs}\pathtag{standard}

A \engterm{system} of $n$ first-order ODEs has the form $\dot{\mathbf{y}}
= \mathbf{f}(t, \mathbf{y})$, where $\mathbf{y} \in \R^{n}$ and
$\mathbf{f} \colon \R \times \R^{n} \to \R^{n}$. The IVP carries an
$n$-vector initial condition $\mathbf{y}(t_{0}) = \mathbf{y}_{0}$.

Every higher-order ODE is a system in disguise. The second-order
equation $\ddot{x} + 2 \zeta \omega_{0} \dot{x} + \omega_{0}^{2} x =
0$ becomes the first-order system
\[
\dot{\mathbf{y}} = \begin{pmatrix} \dot{x} \\ \ddot{x} \end{pmatrix} = \begin{pmatrix} 0 & 1 \\ -\omega_{0}^{2} & -2\zeta\omega_{0} \end{pmatrix} \begin{pmatrix} x \\ \dot{x} \end{pmatrix},
\]
with state vector $\mathbf{y} = (x, \dot{x})^{T}$. The general
$n$th-order ODE becomes an $n$-dimensional first-order system by
the same trick, introducing one new variable for each derivative.
This reduction is why almost every numerical ODE solver is
written for first-order systems.

\subsection{Linear systems with constant coefficients}

A \engterm{linear} system has $\mathbf{f}(t, \mathbf{y}) =
A(t)\mathbf{y} + \mathbf{q}(t)$ with matrix-valued coefficient
$A(t)$ and vector-valued forcing $\mathbf{q}(t)$. With $A$
constant and $\mathbf{q} = 0$, the homogeneous system $\dot{\mathbf{y}}
= A\mathbf{y}$ has solution
\[
\mathbf{y}(t) = e^{A t}\, \mathbf{y}(0),
\]
where $e^{A t}$ is the \engterm{matrix exponential}, defined by
the convergent power series
\[
e^{A t} = \sum_{n = 0}^{\infty} \frac{(A t)^{n}}{n!} = I + A t + \frac{(A t)^{2}}{2!} + \cdots.
\]
The series converges for every square matrix $A$ and every $t$.

For diagonal $A$ with entries $\lambda_{1}, \ldots, \lambda_{n}$,
the matrix exponential is diagonal with entries $e^{\lambda_{1} t},
\ldots, e^{\lambda_{n} t}$. For diagonalisable $A = P D P^{-1}$,
the matrix exponential is $e^{A t} = P e^{D t} P^{-1}$. For
non-diagonalisable $A$ (defective matrices with Jordan blocks),
$e^{A t}$ contains polynomial-times-exponential terms; we leave
the working theory of the Jordan form to Chapter 10.

The eigenvalues of $A$ are the system's \engterm{modes}: an
eigenvalue with negative real part contributes a decaying mode, an
eigenvalue with positive real part a growing mode, an eigenvalue
on the imaginary axis a sustained oscillation, an eigenvalue at the
origin a constant mode. The engineer reads stability off the
spectrum of $A$ at a glance.

\subsection{Computing the matrix exponential in practice}

The series definition is rarely the right way to compute $e^{A t}$
numerically. Three working methods cover most cases.

\textbf{Diagonalisation}, when $A = P D P^{-1}$ with $D$ diagonal,
gives $e^{A t} = P e^{D t} P^{-1}$ exactly, with $e^{D t}$ a
diagonal matrix of scalar exponentials. This is the cleanest path
when $A$ has a complete set of eigenvectors and the eigendecomposition
is well-conditioned (no near-defective spectrum). For a $2 \times
2$ matrix with distinct real eigenvalues, the computation reduces
to two scalar exponentials and one $2 \times 2$ matrix inversion.

\textbf{Scaling and squaring}, the workhorse of production
libraries, uses the identity $e^{A} = (e^{A/2^{s}})^{2^{s}}$ to
reduce the norm of the argument before evaluating a Pad\'e
approximant. The chosen $s$ ensures $\|A/2^{s}\|$ is small enough
for a fixed-order Pad\'e approximation to be accurate to working
precision, then the result is squared back up. The combination is
backward-stable and is the algorithm behind MATLAB's \verb|expm|
and SciPy's \verb|scipy.linalg.expm|.

\textbf{The Jordan form} handles defective matrices (those with
fewer eigenvectors than dimensions) by adding nilpotent blocks
that produce polynomial-times-exponential terms. A Jordan block of
size $k$ with eigenvalue $\lambda$ contributes terms of the form
$e^{\lambda t}, t e^{\lambda t}, \ldots, t^{k-1} e^{\lambda t}/(k-1)!$
to the matrix exponential. Defective matrices arise generically in
critically-damped oscillators and in degenerate parameter settings.
The working theory of the Jordan form lives in Chapter 10; in this
chapter the working knowledge is that polynomial-times-exponential
terms in the solution signal a repeated root of the characteristic
polynomial and a non-diagonalisable Jacobian.

For a $2 \times 2$ system with constant coefficients, the engineer's
closed-form working knowledge of $e^{A t}$ is captured by the
trace--determinant classification: with $T = \operatorname{tr} A$
and $D = \det A$, the eigenvalues are
$\lambda_{\pm} = (T \pm \sqrt{T^{2} - 4 D})/2$, and the four
qualitative regimes (saddle, node, spiral, centre) follow from the
signs of $T$ and the discriminant $T^{2} - 4 D$.

\subsection{Forced linear systems and the variation-of-parameters formula}

The driven system $\dot{\mathbf{y}} = A \mathbf{y} + \mathbf{q}(t)$
with $\mathbf{y}(0) = \mathbf{y}_{0}$ has solution
\[
\mathbf{y}(t) = e^{A t}\, \mathbf{y}_{0} + \int_{0}^{t} e^{A (t - s)}\, \mathbf{q}(s)\,ds.
\]
The first term is the unforced response; the second is the
convolution of the forcing with the system's impulse response. The
integral form is the multivariable analogue of the integrating-factor
formula in §7.2.

\subsection{The heat equation as a system: bridge to PDEs}

Discretising the heat equation $u_{t} = \alpha u_{xx}$ on $N$ grid
points along an interval, with second-difference approximation for
$u_{xx}$, produces a linear system $\dot{\mathbf{u}} = \alpha A
\mathbf{u}$ where $A$ is the tridiagonal matrix with $-2$ on the
diagonal and $1$ on the sub- and super-diagonals (divided by
$\Delta x^{2}$). The eigenvalues of $A$ are real, negative, and
range from approximately $-(\pi/L)^{2}$ for the smoothest mode to
approximately $-4/(\Delta x)^{2}$ for the noisiest mode. The mode
with eigenvalue $\lambda_{j}$ decays with time constant $-1/(\alpha
\lambda_{j})$. This is the transport archetype's first appearance
in matrix form; the apparatus carries through to Chapter 12 (PDEs)
where the continuous version is treated, and to the formal stability
analysis of Chapter 16 (Numerical methods). The transport archetype
runs from a single time constant in the RC circuit, through an
$n$-dimensional set of decay times in the discretised heat equation,
to the continuous spectrum of the full PDE.

\subsection{Nonlinear systems: linearisation}

Most engineering systems are nonlinear. The standard approach
linearises around an \engterm{equilibrium} $\mathbf{y}^{*}$ (a point
where $\mathbf{f}(\mathbf{y}^{*}) = 0$): write $\mathbf{y} =
\mathbf{y}^{*} + \mathbf{\delta}$ and Taylor-expand
$\mathbf{f}(\mathbf{y}^{*} + \mathbf{\delta}) =
J(\mathbf{y}^{*})\,\mathbf{\delta} + O(|\mathbf{\delta}|^{2})$,
where
\[
J = \frac{\partial \mathbf{f}}{\partial \mathbf{y}}
\]
is the Jacobian matrix. The linearised system $\dot{\mathbf{\delta}} =
J(\mathbf{y}^{*})\,\mathbf{\delta}$ controls the local behaviour
near the equilibrium when $J$ is hyperbolic (no eigenvalues on the
imaginary axis); the Hartman-Grobman theorem says the linear
behaviour and the nonlinear behaviour are topologically equivalent
in a neighbourhood of the equilibrium in that case. The proof
belongs in a course on dynamical systems; the working consequence
is that an engineer can certify local stability by computing the
eigenvalues of one matrix.

\section{Phase-plane analysis}\pathtag{standard}

For a two-dimensional autonomous system $\dot{x} = f(x, y)$,
$\dot{y} = g(x, y)$, the \engterm{phase plane} is the $(x, y)$
plane. A solution is a curve $(x(t), y(t))$ traced out as $t$
advances. The vector field $(f, g)$ assigns to each point a
direction and speed; the solution at any point follows the local
direction. The picture summarises every solution at once.

\subsection{Equilibria and their classification by linearisation}

An \engterm{equilibrium} is a point where $f = g = 0$. The
linearisation at an equilibrium has Jacobian $J$ a $2 \times 2$
matrix; its trace $T = J_{11} + J_{22}$ and determinant $D =
J_{11} J_{22} - J_{12} J_{21}$ classify the equilibrium:

\begin{itemize}[leftmargin=*]
  \item $D < 0$: \engterm{saddle}. Two real eigenvalues of opposite
    sign. Always unstable; trajectories approach along one
    eigenvector and recede along the other.
  \item $D > 0$, $T < 0$, $T^{2} > 4D$: \engterm{stable node}. Two
    real negative eigenvalues; trajectories approach the equilibrium
    along the eigenvectors.
  \item $D > 0$, $T < 0$, $T^{2} < 4D$: \engterm{stable spiral}.
    Complex conjugate eigenvalues with negative real part;
    trajectories spiral in.
  \item $D > 0$, $T > 0$, $T^{2} > 4D$: \engterm{unstable node}.
    Two real positive eigenvalues; trajectories recede.
  \item $D > 0$, $T > 0$, $T^{2} < 4D$: \engterm{unstable spiral}.
    Complex eigenvalues with positive real part.
  \item $D > 0$, $T = 0$: \engterm{centre}. Pure imaginary
    eigenvalues; closed orbits in the linearisation. The
    classification is fragile under nonlinear corrections; the
    nonlinear system may have a stable spiral, an unstable spiral,
    or genuine closed orbits (the case for Hamiltonian systems).
\end{itemize}

The trace-determinant plane summarises all four cases on a single
picture: the parabola $T^{2} = 4 D$ separates real from complex
eigenvalues, and the axes $D = 0$ and $T = 0$ separate the
stability regimes. The engineer who has internalised the
trace-determinant picture reads phase-plane behaviour off the
linearisation by inspection.

% Figure: trace-determinant plane for 2x2 linear systems. A full-page
% reference picture.
\begin{figure}[p]
\centering
\begin{tikzpicture}[
  axis/.style={draw=engblack, -{Stealth}, thick},
  region/.style={font=\small, align=center},
  parabola/.style={draw=engblue, thick},
  label/.style={font=\small},
]
% Axes: T (trace) horizontal, D (determinant) vertical
\draw[axis] (-5.5,0) -- (5.5,0) node[anchor=west, label] {$T = \operatorname{tr} J$};
\draw[axis] (0,-3.0) -- (0,5.5) node[anchor=south, label] {$D = \det J$};

% Ticks
\foreach \x in {-4,-2,2,4} {
  \draw (\x,-0.08) -- (\x,0.08);
  \node[font=\scriptsize, below] at (\x,-0.12) {\x};
}
\foreach \y in {-2,2,4} {
  \draw (-0.08,\y) -- (0.08,\y);
  \node[font=\scriptsize, left] at (-0.12,\y) {\y};
}

% Parabola T^2 = 4D, i.e. D = T^2 / 4
\draw[parabola] plot[domain=-5:5, samples=200] ({\x}, {(\x)^2/4});
\node[engblue, font=\scriptsize, anchor=west] at (3.2, 3.5) {$T^{2}=4D$};

% Shade the negative-D region (saddles)
\fill[engred, opacity=0.1] (-5.5,-3) rectangle (5.5,0);
% Shade the right half of upper region (unstable)
\fill[engamber, opacity=0.08] (0,0) -- (5.5,0) -- (5.5,5.5) -- (0,5.5) -- cycle;
% Shade the left half of upper region (stable)
\fill[englime, opacity=0.10] (0,0) -- (-5.5,0) -- (-5.5,5.5) -- (0,5.5) -- cycle;

% Region labels
\node[region, engred] at (0,-1.7) {\textbf{Saddle} \\ ($D<0$, unstable)};
\node[region, engblack] at (-3.3, 4.6) {\textbf{Stable spiral} \\ ($T<0$, $T^{2}<4D$)};
\node[region, engblack] at (-3.3, 1.2) {\textbf{Stable node} \\ ($T<0$, $T^{2}>4D$)};
\node[region, engblack] at (3.3, 4.6) {\textbf{Unstable spiral} \\ ($T>0$, $T^{2}<4D$)};
\node[region, engblack] at (3.3, 1.2) {\textbf{Unstable node} \\ ($T>0$, $T^{2}>4D$)};
\node[region, engblue, anchor=south] at (0, 5.6) {\textbf{Centre} ($T=0$)};

\end{tikzpicture}
\caption[The trace--determinant plane for $2\times 2$ linear
systems]{The trace--determinant plane classifies equilibria of
linear systems $\dot{\mathbf{y}}=J\mathbf{y}$ in two dimensions.
The parabola $T^{2}=4D$ separates real eigenvalues (below) from
complex eigenvalues (above). The axis $D=0$ separates saddles
($D<0$, opposite-sign real eigenvalues) from nodes and spirals
($D>0$). The axis $T=0$ separates stable behaviour ($T<0$, both
eigenvalues have negative real part) from unstable behaviour
($T>0$). The boundary $T=0$, $D>0$ is the centre case, with pure
imaginary eigenvalues and closed orbits in the linearisation; the
classification is fragile under nonlinear corrections, which may
push the equilibrium toward a stable spiral, an unstable spiral,
or genuine closed orbits (in Hamiltonian systems). The
trace--determinant picture is the working diagram an engineer
internalises for phase-plane diagnosis at a glance.}
\label{fig:vol02:ch07:trace-determinant}
\end{figure}


% Figure: phase-plane portrait of the damped harmonic oscillator
% (underdamped), showing the spiral trajectory toward the origin.
\begin{figure}[htbp]
\centering
\begin{tikzpicture}[
  axis/.style={draw=engblack, -{Stealth}, thick},
  field/.style={draw=enggray, thin, -{Latex[length=1.2mm, width=0.9mm]}},
  traj/.style={draw=engblue, very thick},
  equilib/.style={draw=engred, fill=engred, circle, inner sep=1.5pt},
  label/.style={font=\small},
]
% Centered axes: x from -5 to 5 (position), y from -5 to 5 (velocity)
\draw[axis] (-5.2,0) -- (5.2,0) node[anchor=west, label] {$x$};
\draw[axis] (0,-5.2) -- (0,5.2) node[anchor=south, label] {$\dot{x}$};

\foreach \x in {-4,-2,2,4} {
  \draw (\x,-0.07) -- (\x,0.07);
  \node[font=\scriptsize, below] at (\x,-0.1) {\x};
}
\foreach \y in {-4,-2,2,4} {
  \draw (-0.07,\y) -- (0.07,\y);
  \node[font=\scriptsize, left] at (-0.1,\y) {\y};
}

% Vector field for damped oscillator: xdot = y, ydot = -omega0^2 * x - 2*zeta*omega0 * y
% omega0 = 1, zeta = 0.10
\foreach \x in {-4.5,-3,-1.5,0,1.5,3,4.5} {
  \foreach \y in {-4.5,-3,-1.5,0,1.5,3,4.5} {
    \pgfmathsetmacro{\fx}{\y}
    \pgfmathsetmacro{\fy}{-1*\x - 0.2*\y}
    \pgfmathsetmacro{\norm}{sqrt((\fx)^2 + (\fy)^2)}
    \ifnum 1=1
      \pgfmathsetmacro{\dx}{0.35*\fx/(\norm+0.001)}
      \pgfmathsetmacro{\dy}{0.35*\fy/(\norm+0.001)}
      \draw[field] (\x - \dx, \y - \dy) -- (\x + \dx, \y + \dy);
    \fi
  }
}

% Trajectory: spiral inward from (4, 0). For underdamped osc with omega0=1, zeta=0.10:
%   x(t) = exp(-zeta * t) * (A cos(omega_d t) + B sin(omega_d t))
%   with omega_d = sqrt(1 - zeta^2) ~ 0.995, zeta = 0.10
%   Initial: x(0) = 4, xdot(0) = 0 => A = 4, B = 4*zeta/omega_d ~ 0.402
\draw[traj] plot[domain=0:30, samples=400] (
  {exp(-0.10*\x) * (4 * cos(\x * 0.995 r) + 0.402 * sin(\x * 0.995 r))},
  {exp(-0.10*\x) * (-4 * sin(\x * 0.995 r) * 0.995 + 0.402 * cos(\x * 0.995 r) * 0.995 - 0.10 * (4 * cos(\x * 0.995 r) + 0.402 * sin(\x * 0.995 r)))}
);

% Equilibrium dot
\node[equilib] at (0,0) {};
\node[font=\scriptsize, engred, anchor=south west] at (0.1, 0.1) {stable spiral};

% Initial-condition marker
\node[font=\scriptsize, engblue, anchor=west] at (4.1, 0) {$\mathbf{y}_{0}$};
\draw[engblue, fill=engblue] (4,0) circle (1.2pt);

\end{tikzpicture}
\caption[Phase portrait of the underdamped oscillator]{Phase
portrait of the damped harmonic oscillator $\ddot{x}+2\zeta\omega_{0}
\dot{x}+\omega_{0}^{2}x=0$ with $\omega_{0}=1$ and $\zeta=0.10$.
The vector field $(\dot{x},\ddot{x})=(\dot{x},-\omega_{0}^{2}x-2\zeta
\omega_{0}\dot{x})$ assigns a direction at every point; the
trajectory through $(x_{0},\dot{x}_{0})=(4,0)$ spirals inward to the
origin. The origin is a \emph{stable spiral}, classified by the
linearisation's complex eigenvalues with negative real part. For
$\zeta=0$ the spiral becomes a closed orbit (centre); for $\zeta>1$
it becomes a stable node (no rotation).}
\label{fig:vol02:ch07:phase-portrait}
\end{figure}


\subsection{Worked example: the predator-prey system}

The Lotka-Volterra equations are
\[
\dot{x} = a x - b x y, \quad \dot{y} = -c y + d x y,
\]
with $a, b, c, d > 0$. The state $x$ is the prey population, $y$
the predator population. There are two equilibria: $(0, 0)$ and
$(c/d, a/b)$.

At $(0, 0)$: the Jacobian is $\operatorname{diag}(a, -c)$. Trace
$T = a - c$, determinant $D = -a c < 0$. Saddle: prey grows in
absence of predators, predators decay in absence of prey.

At $(c/d, a/b)$: the Jacobian is
\[
J = \begin{pmatrix} 0 & -b c/d \\ a d/b & 0 \end{pmatrix},
\]
with trace $T = 0$ and determinant $D = a c > 0$. The linearisation
predicts a centre. The full nonlinear system has closed orbits
(conserved quantity: $H(x, y) = d x - c \ln x + b y - a \ln y$;
see the exercises). Predator and prey populations oscillate, with
amplitude set by the initial condition.

The model is a textbook idealisation; real ecology adds finite
carrying capacity, age structure, refuge effects, and stochastic
forcing. The lesson the model teaches at first reading is the
phase-plane picture: a centre with closed orbits, oscillations
without damping. The lesson on second reading is that the centre
is fragile: small modelling additions push the equilibrium toward
a stable spiral (typical) or an unstable spiral.

\subsection{Lyapunov stability, previewed}

A direct certificate of stability, when one is available, comes
from a \engterm{Lyapunov function}: a scalar function $V(\mathbf{y})$
that is positive in a neighbourhood of the equilibrium except at
the equilibrium itself, where it vanishes, and whose time
derivative along trajectories is non-positive. If $V$ exists, the
equilibrium is stable; if $\dot{V} < 0$ strictly off the
equilibrium, the equilibrium is asymptotically stable. The
function $V$ is often a physical energy: kinetic plus potential for
a mechanical oscillator, electrical plus magnetic for an LC
circuit. The full Lyapunov apparatus, including the construction
of $V$ for systems where physical energy is not the right
candidate, is in Volume X; we mention it here so the reader can
see the picture.

\section{Numerical ODE solvers: Euler, RK4, stiffness}\pathtag{core}

Most ODEs encountered in working engineering have no closed-form
solution. The engineer integrates them numerically. The chapter
covers three solver classes (forward Euler, classical RK4, implicit
Euler) and the working diagnosis of \engterm{stiffness}. The full
catalogue of methods, the formal order of accuracy proofs, and the
modern adaptive-step solvers belong in Chapter 16.

\subsection{Forward Euler}

The simplest method advances $\mathbf{y}_{n+1} = \mathbf{y}_{n} + h
\mathbf{f}(t_{n}, \mathbf{y}_{n})$, with $h$ the \engterm{step
size} and $\mathbf{y}_{n} \approx \mathbf{y}(t_{n})$. The
right-hand side at the current point is the slope estimate; the
step extrapolates linearly. The local truncation error per step is
$O(h^{2})$, and the global error after $N = T/h$ steps is $O(h)$:
forward Euler is a \engterm{first-order} method.

\subsection{Classical fourth-order Runge-Kutta (RK4)}

RK4 advances by combining four slope estimates at the start, two
intermediate points, and the end of each step:
\[
\begin{aligned}
\mathbf{k}_{1} &= \mathbf{f}(t_{n}, \mathbf{y}_{n}), \\
\mathbf{k}_{2} &= \mathbf{f}(t_{n} + h/2, \mathbf{y}_{n} + h\mathbf{k}_{1}/2), \\
\mathbf{k}_{3} &= \mathbf{f}(t_{n} + h/2, \mathbf{y}_{n} + h\mathbf{k}_{2}/2), \\
\mathbf{k}_{4} &= \mathbf{f}(t_{n} + h, \mathbf{y}_{n} + h\mathbf{k}_{3}), \\
\mathbf{y}_{n+1} &= \mathbf{y}_{n} + \frac{h}{6}(\mathbf{k}_{1} + 2\mathbf{k}_{2} + 2\mathbf{k}_{3} + \mathbf{k}_{4}).
\end{aligned}
\]
The method's local truncation error is $O(h^{5})$, and its global
error is $O(h^{4})$: RK4 is a \engterm{fourth-order} method. For
the same step size, RK4 is roughly four right-hand-side
evaluations per step, four times the cost of Euler, but the error
falls as $h^{4}$ instead of $h$, so a tenfold reduction in step
size buys $10^{4}$ times the accuracy with $40\times$ the work.
For smooth, non-stiff right-hand sides at moderate accuracy, RK4 is
the standard first-pass choice.

\subsection{Implicit Euler and the stability story}

The \engterm{implicit Euler} method solves at each step
$\mathbf{y}_{n+1} = \mathbf{y}_{n} + h \mathbf{f}(t_{n+1},
\mathbf{y}_{n+1})$. The new state appears on both sides;
$\mathbf{y}_{n+1}$ is computed by solving a (possibly nonlinear)
algebraic equation per step. The method is also first-order
accurate, but it has a property forward Euler lacks: it is
unconditionally stable for the test equation $\dot{y} = \lambda y$
with $\lambda < 0$.

The test equation $\dot{y} = \lambda y$ exposes the stability
distinction. Forward Euler gives $y_{n+1} = (1 + h \lambda) y_{n}$;
the iterates remain bounded only if $|1 + h \lambda| \leq 1$,
which for $\lambda$ real negative requires $h \leq 2/|\lambda|$.
Implicit Euler gives $y_{n+1} = y_{n}/(1 - h \lambda)$; for
$\lambda$ real negative, $|1 - h \lambda| \geq 1$ always, so the
iterates always damp. The cost is the per-step solve.

\subsection{Stiffness}

A system is \engterm{stiff} when its eigenvalues span a wide range
of magnitudes, the fastest decaying mode is much faster than the
slowest, and the engineer cares about the long-term behaviour.
Forward Euler must use a step size set by the fastest mode
(stability constraint) even though the slow modes are what the
engineer wants to track; the resulting computation is dominated by
short-step integration of fast modes that have already decayed.
Implicit methods integrate the fast modes correctly without the
short-step penalty.

The classic test: $\dot{y}_{1} = -1000 y_{1}$, $\dot{y}_{2} =
-y_{2}$. The fast mode decays in milliseconds; the slow mode in
seconds. Forward Euler requires $h < 2/1000 = 2\,\si{\milli\second}$
for stability even after the fast mode has already decayed below
machine precision. Implicit Euler can use $h = 0.1\,\si{\second}$,
$50\times$ larger, and produce the slow mode accurately.

% Figure: forward Euler stability region in the complex h*lambda plane.
% Compares explicit and implicit Euler.
\begin{figure}[htbp]
\centering
\begin{tikzpicture}[
  axis/.style={draw=engblack, -{Stealth}, thick},
  unstable/.style={draw=engred, very thick, fill=engred, fill opacity=0.10},
  stable/.style={draw=engblue, very thick, fill=engblue, fill opacity=0.10},
  label/.style={font=\small},
]
% Two side-by-side panels: explicit Euler (left), implicit Euler (right)
% Panel 1: explicit Euler stability region |1 + z| <= 1, circle around (-1, 0) of radius 1
\begin{scope}[xshift=0cm]
  \draw[axis] (-3.5,0) -- (1.5,0) node[anchor=west, label] {$\operatorname{Re}(h\lambda)$};
  \draw[axis] (0,-2.2) -- (0,2.2) node[anchor=south, label] {$\operatorname{Im}(h\lambda)$};
  \foreach \x in {-3,-2,-1,1} {
    \draw (\x,-0.07) -- (\x,0.07);
    \node[font=\scriptsize, below] at (\x,-0.1) {\x};
  }
  \foreach \y in {-2,-1,1,2} {
    \draw (-0.07,\y) -- (0.07,\y);
    \node[font=\scriptsize, left] at (-0.1,\y) {\y};
  }
  \fill[engblue, fill opacity=0.20, draw=engblue, very thick] (-1,0) circle (1.0);
  \node[engblue, font=\small] at (-1, 0) {stable};
  \node[engblack, font=\small\bfseries, anchor=north] at (-1, -2.3) {Forward Euler};
\end{scope}

% Panel 2: implicit Euler stability region |1 - z| >= 1, exterior of unit
% disc around (+1, 0)
\begin{scope}[xshift=8cm]
  \draw[axis] (-3.5,0) -- (1.5,0) node[anchor=west, label] {$\operatorname{Re}(h\lambda)$};
  \draw[axis] (0,-2.2) -- (0,2.2) node[anchor=south, label] {$\operatorname{Im}(h\lambda)$};
  \foreach \x in {-3,-2,-1,1} {
    \draw (\x,-0.07) -- (\x,0.07);
    \node[font=\scriptsize, below] at (\x,-0.1) {\x};
  }
  \foreach \y in {-2,-1,1,2} {
    \draw (-0.07,\y) -- (0.07,\y);
    \node[font=\scriptsize, left] at (-0.1,\y) {\y};
  }
  % Shade the whole rectangle blue, then erase the unit disc around (1, 0)
  \begin{scope}
    \clip (-3.5,-2.2) rectangle (1.5, 2.2);
    \fill[engblue, fill opacity=0.20] (-3.5,-2.2) rectangle (1.5, 2.2);
    \fill[white] (1,0) circle (1.0);
  \end{scope}
  \draw[engblue, very thick] (1,0) circle (1.0);
  \node[engblue, font=\small] at (-2.0, 1.4) {stable};
  \node[engred, font=\small] at (1, 0) {unstable};
  \node[engblack, font=\small\bfseries, anchor=north] at (-1, -2.3) {Implicit Euler};
\end{scope}

\end{tikzpicture}
\caption[Stability regions for forward and implicit
Euler]{Stability regions in the complex $h\lambda$ plane for
forward Euler (left) and implicit Euler (right) applied to the
test equation $\dot{y}=\lambda y$. Forward Euler's stability
region is the disc $|1+h\lambda|\leq 1$, a circle of radius one
centred at $-1$; the method is stable only for $h\lambda$ inside
this disc. Implicit Euler's stability region is the exterior of
the disc $|1-h\lambda|\leq 1$ centred at $+1$, which includes the
entire left half-plane. A method is \emph{A-stable} when its
stability region contains the left half-plane; implicit Euler is
A-stable, forward Euler is not. Stiffness is the condition under
which the step size demanded by accuracy is much larger than the
step size demanded by forward-Euler stability; A-stable methods
escape the stability constraint and pay the per-step cost of an
implicit solve.}
\label{fig:vol02:ch07:stiffness}
\end{figure}


The diagnosis of stiffness is partly mathematical (compute the
Jacobian's spectrum) and partly empirical (the explicit solver
takes thousands of steps to advance one slow time scale). Modern
adaptive solvers detect stiffness by monitoring step-size
rejection, but the engineer must know the symptom and the response.

\subsection{Step-size error analysis: a worked example}

Take $\dot{y} = -y$, $y(0) = 1$, exact solution $y(t) = e^{-t}$.
At $t = 1$, the exact value is $e^{-1} \approx 0.3679$.

Forward Euler with $h = 0.1$ requires ten steps. After ten steps,
$y_{10} = (1 - 0.1)^{10} = 0.9^{10} \approx 0.3487$. The error is
$0.0192$, about $5.2\,\%$. Halving the step size to $h = 0.05$
gives $y_{20} = (0.95)^{20} \approx 0.3585$, error $0.0094$,
about $2.6\,\%$. The error halves with the step size: first-order.

RK4 with $h = 0.1$ gives $y_{10} \approx 0.367879$; the error is
about $4 \times 10^{-7}$. Halving the step size to $h = 0.05$
reduces the error by a factor of $16$ (fourth-order), well below
single-precision representation. For most non-stiff problems at
working engineering tolerances, RK4 is the default; the step size
chosen by accuracy lies far above the stability bound.

\section{Failure: numerical instability and the stiffness trap}\pathtag{core}

ODE solvers fail in three distinguishable ways: by numerical
instability (the answer grows because the method does, not because
the system does), by slow accumulation of arithmetic error (the
answer drifts because each step rounds), and by silent migration
between attractors (the answer is qualitatively wrong because the
method's modified equation has different long-term behaviour).
Each failure has a working diagnosis. Each has an institutional
example.

\subsection{Numerical instability}

The forward-Euler stability constraint $h \leq 2/|\lambda|$ for the
linear test equation generalises (in approximate form) to any
linear system: the step size must be smaller than $2/\rho(A)$,
where $\rho(A)$ is the spectral radius of the Jacobian. For a
stiff system whose largest eigenvalue is $-10^{6}\,\si{\per
\second}$, the explicit solver requires $h < 2 \times 10^{-6}\,\si
{\second}$, even if the engineer cares about a one-second
phenomenon. Running an explicit solver on a stiff problem with too
large a step does not produce a slowly diverging answer; it
produces an exploding answer that the user, watching the screen,
typically attributes to a bug in the right-hand side rather than to
a stability violation.

The diagnosis: the answer oscillates with growing amplitude,
typically with period $2 h$ (the Nyquist frequency of the
discretisation). The response: switch to an implicit method, or
shrink the step size below the stability bound. Both work; only
the first is cost-effective on stiff problems.

\subsection{The stiffness trap}

Stiffness traps the unwary in two ways. First, an engineer who
reaches for a default explicit solver on a stiff problem watches
the integration take thousands of steps per unit time and
typically concludes either that the problem is hard or that the
solver is broken; the correct conclusion is that the method is
mismatched to the problem. Second, a problem that starts non-stiff
can become stiff as a slow mode catches up to a fast one or as a
limit cycle approaches a relaxation regime; an adaptive solver
that does not detect the transition will silently switch from
accurate-but-slow to fast-but-wrong, depending on which heuristic
its step controller follows.

The standard discipline:

\begin{enumerate}[leftmargin=*]
  \item Estimate the spectral range of the right-hand side's
    Jacobian. If the ratio of largest to smallest magnitude exceeds
    $10$ to $100$, suspect stiffness.
  \item Run the candidate explicit solver at modest tolerance and
    inspect the per-step cost. Thousands of steps per unit time on
    a smooth problem signals stiffness.
  \item Switch to an implicit method (BDF, implicit Runge-Kutta,
    Rosenbrock) when stiffness is suspected. The cost per step
    rises (a Newton solve is required); the cost per unit time
    typically falls.
  \item Verify the integrated solution against an analytical or
    asymptotic answer in the regime where one exists. A solver that
    silently produces wrong answers in the easy regime will not be
    rescued by the hard regime.
\end{enumerate}

\subsection{Accumulated arithmetic error: the Patriot pattern in numerical integration}

The Patriot missile system at Dhahran (1991) is the canonical
software case of accumulated arithmetic error in a long-running
clock; the registry entry is at
\repopath{docs/research/accidents/patriot-dhahran-1991.md}, and the
same mechanism appears in long-running numerical ODE integration in
a form the ODE engineer must recognise on sight. The case warrants
unpacking at three scales: the technical mechanism, the
organisational context, and the lesson for numerical integration.

\paragraph{Technical mechanism.} On 25 February 1991, a US Army
Patriot missile battery at Dhahran, Saudi Arabia, failed to engage
an incoming Iraqi Scud missile that struck a US Army barracks. The
incident killed 28 service members and wounded approximately 100
others. The U.S.\ General Accounting Office investigated under
congressional request and issued report GAO/IMTEC-92-26 in
February 1992 \cite{acc:gao-patriot-1992}. The GAO identified the
cause as a software-clock arithmetic error in the Patriot's
range-gate prediction.

The Patriot's tracking software represented elapsed time as the
number of tenths of a second since the system was started, stored
as an integer counter, and converted to seconds as a real number
when needed. The conversion used a $24$-bit fixed-point
representation of the constant $1/10$. The exact decimal $0.1$ has
no exact binary representation; the $24$-bit truncation introduced
a small per-conversion error of approximately $9.5 \times 10^{-8}$
in the converted seconds value. The error accumulated linearly with
the integer counter.

After approximately $100\,\si{\hour}$ of continuous operation, the
integer counter had reached approximately $3.6 \times 10^{6}$, and
the cumulative conversion error had reached approximately
$0.34\,\si{\second}$ \cite{acc:gao-patriot-1992}. The Patriot's
range-gate calculation used the converted time to predict where an
incoming target would be at the next radar look. At the Scud's
speed of approximately $1{,}676\,\si{\meter\per\second}$, the
$0.34\,\si{\second}$ time offset corresponded to a predicted-position
error of approximately $570\,\si{\meter}$, larger than the range
gate. The system reported "no track" and did not engage.

\paragraph{Organisational context.} The GAO report identified
contributing institutional factors beyond the technical error
\cite{acc:gao-patriot-1992}. Patriot batteries had not been
designed to operate continuously for one hundred hours; the
operational concept assumed periodic restarts that would reset the
time counter. Field operators in the Gulf War deployment had been
instructed to keep the systems running continuously to maintain
readiness. The accumulated-error effect had been observed in field
tests and reported through the chain. A software patch correcting
the conversion error had been developed and was being distributed
to Patriot batteries at the time of the Dhahran incident; the patch
arrived at the Dhahran battery on 26 February 1991, the day after
the incident. The patch and the operational-restart guidance had
moved through different distribution channels with no
synchronisation; a battery in the field could have one without the
other. This pattern, in which a technical correction and an
operational workaround propagate through separate institutional
paths and arrive at field units inconsistently, recurs across
software-mediated systems and is the organisational analogue of the
arithmetic mechanism.

\paragraph{Lesson for ODE modelling.} The mechanism transfers
directly to numerical ODE integration. Each step's roundoff error
is bounded above by machine epsilon times the working numbers, but
errors accumulate. For a non-stiff stable solver running for $N$
steps with floating-point machine epsilon $\epsilon$ and per-step
propagation factor close to one, the cumulative error grows roughly
as $\sqrt{N}\, \epsilon\, \|y\|$ (random-walk roundoff, when
per-step rounding directions are effectively independent) or as
$N \epsilon \|y\|$ (worst case, correlated roundoff). The Patriot's
accumulated error is the latter regime: $3.6 \times 10^{6}$ counter
increments, each rounded by about $9.5 \times 10^{-8}$ in a
consistent direction, give the linear product $3.6 \times 10^{6}
\times 9.5 \times 10^{-8} \approx 0.34$, matching the GAO figure.

A naive RK4 integration of a smooth oscillator at IEEE double
precision ($\epsilon \approx 2.2 \times 10^{-16}$) for $N = 10^{9}$
steps accumulates worst-case roundoff of order $10^{9} \times 2.2
\times 10^{-16} \approx 2 \times 10^{-7}$, comparable to the
per-step truncation error at $h = 10^{-3}$ on a tolerant problem.
For longer runs (a multi-year orbital simulation, a multi-decade
climate run, a continuous control loop), the run-length budget
must be set explicitly, the accumulated-roundoff estimate tracked,
and the precision raised or the time horizon shortened when the
budget is exceeded. The discipline is to know the run length, the
per-step error, and the application's accuracy budget before
starting. Modern simulation codes for climate, plasma, fluid, and
orbital mechanics maintain runtime estimates of accumulated
roundoff and switch to extended precision when the accumulated
estimate crosses a tolerance. The discipline is the numerical-
integration counterpart of the operational-restart policy that the
Patriot's design had assumed and the Dhahran deployment had
discarded.

\subsection{Reuse outside the qualification envelope: Ariane 5 Flight 501}

The Ariane 5 maiden flight on 4 June 1996 broke up about $39\,\si{
\second}$ after main engine ignition when an inertial reference
system software exception, triggered by a $64$-bit-to-$16$-bit
conversion overflow inside an alignment function reused unchanged
from Ariane 4, was unhandled by both the primary and the back-up
SRI; the launcher then commanded extreme manoeuvres on bad
attitude data and broke up under aerodynamic loads
\cite{acc:ariane501-ifr}. The registry entry
\repopath{docs/research/accidents/ariane-501-1996.md} carries the
full account. For an Ariane 4 trajectory, the BH (horizontal
bias) variable that overflowed had been within $16$-bit signed
integer range; for the Ariane 5 trajectory, which built horizontal
velocity faster, it was not. The variable had been one of seven
that the Ariane 4 software protected against overflow; three of
the others, including BH, had been left unprotected on the basis
that the Ariane 4 trajectory could not produce out-of-range
values. That analysis was not redone for Ariane 5.

For ODE integration, the lesson is the range-versus-envelope
failure mode in numerical form: a solver, a step controller, or a
right-hand-side implementation qualified for one regime is reused
in another whose dynamic range has not been re-checked. A pressure
calculation that worked on subsonic flows overflows in transonic
ones; an attitude propagator validated for short hops produces
wrong angles for long ones; a chemical-kinetics solver tuned for
mild stiffness silently fails when a new reaction step pushes the
stiffness ratio past the integrator's stability region. The
discipline is to treat reused numerical code as new code whose
operating envelope must be re-derived from the new application,
and to run the range analysis (interval arithmetic, fuzzed input
sweeps, formal verification of bounds where the cost justifies it)
against the new envelope before the code carries flight-critical
load.

\subsection{Silent attractor migration}

A modified-equation analysis reveals that a numerical scheme does
not solve the original ODE exactly; it solves a modified ODE that
matches the original to the order of the method's accuracy, plus
higher-order terms. For a Hamiltonian system, the modified
equation generally is not Hamiltonian, and the modified system's
trajectories generally do not conserve the original's energy. A
naive RK4 integration of the predator-prey centre over many
oscillations either spirals slowly inward or slowly outward,
because RK4's modified equation has weak damping or anti-damping.
The system does not have damping; the integrator does. The
solution to a Hamiltonian problem with naive RK4 will eventually
collapse to the equilibrium or escape, neither of which is a
property of the original.

The response is \engterm{symplectic} integrators (leapfrog,
Verlet, symplectic Runge-Kutta), which are designed so the
modified equation remains Hamiltonian, and the long-term energy
behaviour of the integration matches the long-term energy
behaviour of the system. Orbital mechanics, molecular dynamics,
and plasma simulation use symplectic methods routinely. The full
working theory is in Chapter 16.

\section{Worked examples}\pathtag{standard}

We close with four worked applications across domains, each
following the same pattern: dimensional analysis, model setup,
solution (analytical, numerical, or both), and a note on what
fails or accumulates.

\subsection{The RC discharge in detail}

The discharge $RC \dot{v} + v = 0$, $v(0) = V_{0}$, has solution
$v(t) = V_{0} e^{-t/\tau}$ with $\tau = RC$. Quick numerical
check: an RC circuit with $R = 10\,\si{\kilo\ohm}$, $C =
100\,\si{\micro\farad}$, $V_{0} = 5\,\si{\volt}$ has $\tau = 1\,\si
{\second}$, $v(1) = 5/e \approx 1.84\,\si{\volt}$, $v(5)
\approx 0.034\,\si{\volt}$.

What fails: the linear discharge model assumes a constant capacitance.
Real electrolytic capacitors lose capacitance with age, voltage
swing, and temperature; ceramic capacitors with class-2 dielectrics
lose capacitance with applied voltage. A precision discharge
measurement requires the capacitor type and recent calibration
state in the laboratory record. The first-order model is not the
answer; it is the answer near the working point of a freshly
calibrated capacitor.

\subsection{The simple pendulum}

A point mass $m$ at the end of a rigid massless rod of length
$\ell$ in a gravitational field $g$, with no friction, swinging in
a vertical plane, satisfies
\[
\ddot{\theta} + \frac{g}{\ell} \sin\theta = 0,
\]
with $\theta$ the angle from the downward vertical. The equation
is nonlinear; the small-angle approximation $\sin\theta \approx
\theta$ produces the linear oscillator $\ddot{\theta} +
\omega_{0}^{2} \theta = 0$ with $\omega_{0} = \sqrt{g/\ell}$, period
$T_{0} = 2\pi\sqrt{\ell/g}$.

The first amplitude correction comes from expanding $\sin\theta =
\theta - \theta^{3}/6 + O(\theta^{5})$ and applying perturbation
methods (the working derivation is in §7.10's exercise);
\[
T(\theta_{0}) \approx T_{0}\!\left(1 + \frac{\theta_{0}^{2}}{16} + \frac{11 \theta_{0}^{4}}{3072} + \cdots\right).
\]
For $\theta_{0} = 10^{\circ} \approx 0.175$ rad, the first
correction is $0.175^{2}/16 \approx 1.9 \times 10^{-3}$, or
$0.19\,\%$. For $\theta_{0} = 90^{\circ}$ the first correction
is $\pi^{2}/64 \approx 0.154$, or $15\,\%$, and higher-order
terms contribute several percent more; the small-angle formula
under-predicts the period materially. The chapter project
quantifies the error across the full range.

What fails: the small-angle approximation drops a $\theta^{3}$
term that becomes comparable to $\theta$ at large amplitudes. The
linearised oscillator is exact for an oscillator the engineer has
designed to be linear (sufficiently stiff spring, sufficiently
small driver); it is approximate for a pendulum operated near its
linearity boundary; it is wrong for a pendulum swung past
$45^{\circ}$ if the engineer wants better than five-percent
period accuracy.

% Figure: pendulum period as a function of amplitude. Compares the
% small-angle prediction, the first-correction prediction, and the
% exact (elliptic-integral) result.
\begin{figure}[htbp]
\centering
\begin{tikzpicture}[
  axis/.style={draw=engblack, -{Stealth}, thick},
  small/.style={draw=engblue, very thick, dashed},
  first/.style={draw=engamber, very thick, dash dot},
  exact/.style={draw=engred, very thick},
  guide/.style={draw=enggray, dashed, thin},
  label/.style={font=\small},
]
% Axes: x amplitude in degrees (0 to 120, plotted x in 0..10), y T/T0 (1.0 to 1.4, plotted 0..5)
\draw[axis] (0,0) -- (10.5,0) node[anchor=west, label] {$\theta_{0}$ [deg]};
\draw[axis] (0,0) -- (0,5.5)  node[anchor=south, label] {$T/T_{0}$};

% x ticks (12 deg per unit)
\foreach \x/\m in {0/0, 2/24, 4/48, 6/72, 8/96, 10/120} {
  \draw (\x,-0.08) -- (\x,0.08);
  \node[font=\scriptsize, below] at (\x,-0.1) {\m};
}
% y ticks (0.08 per unit; 1.0 at y=0; 1.40 at y=5)
\foreach \y/\h in {0/1.00, 1.25/1.10, 2.5/1.20, 3.75/1.30, 5/1.40} {
  \draw (-0.08,\y) -- (0.08,\y);
  \node[font=\scriptsize, left] at (-0.1,\y) {\h};
}

% Reference line at T/T0 = 1.0
\draw[guide] (0, 0) -- (10, 0);
% Reference line at T/T0 = 1.05 (y = 0.625)
\draw[guide] (0, 0.625) -- (10, 0.625);
\node[engblack, font=\scriptsize, anchor=west] at (10.05, 0.625) {$+5\,\%$};

% Small-angle prediction: T/T0 = 1 always.
\draw[small] (0, 0) -- (10, 0);
\node[engblue, font=\scriptsize, anchor=west] at (10.05, 0) {small-angle};

% First correction: T/T0 = 1 + theta^2/16, theta in radians
% x = theta_deg / 12, theta_deg = 12*x, theta_rad = 12*x*pi/180 = x*pi/15
% Correction = (x*pi/15)^2 / 16; T/T0 - 1 plotted as 12.5 * correction
\draw[first] plot[domain=0:10, samples=120]
  ({\x}, {12.5 * ((\x * 3.14159/15)^2 / 16)});
\node[engamber, font=\scriptsize, anchor=west] at (8.0, 1.7) {first correction $1 + \theta_{0}^{2}/16$};

% Exact (elliptic-integral) values: T/T0 = (2/pi) K(sin(theta/2))
% Sampled at theta_deg = 0, 10, 20, ..., 120
% K(k) values: theta/2 in degrees, k = sin(theta/2 in rad)
% Reference values (Abramowitz & Stegun):
% theta=0:    T/T0 = 1.00000
% theta=10:   T/T0 = 1.00191
% theta=20:   T/T0 = 1.00767
% theta=30:   T/T0 = 1.01741
% theta=40:   T/T0 = 1.03135
% theta=50:   T/T0 = 1.04978
% theta=60:   T/T0 = 1.07314
% theta=70:   T/T0 = 1.10200
% theta=80:   T/T0 = 1.13715
% theta=90:   T/T0 = 1.17984
% theta=100:  T/T0 = 1.23198
% theta=110:  T/T0 = 1.29661
% theta=120:  T/T0 = 1.37864
% Translate: x = theta_deg/12, y = 12.5 * (T/T0 - 1)
\draw[exact] plot coordinates {
  (0,      0.000)
  (0.833,  0.024)
  (1.667,  0.096)
  (2.500,  0.218)
  (3.333,  0.392)
  (4.167,  0.622)
  (5.000,  0.914)
  (5.833,  1.275)
  (6.667,  1.714)
  (7.500,  2.248)
  (8.333,  2.900)
  (9.167,  3.708)
  (10.000, 4.733)
};
\node[engred, font=\small, anchor=south west] at (5.5, 3.4) {exact (elliptic)};

% Mark amplitude at which exact crosses 5% (~67 deg). y=0.625 at x ~ 5.58
\draw[guide] (5.58, 0) -- (5.58, 0.625);
\node[font=\scriptsize, enggray, anchor=south west] at (5.5, 0.0) {$\sim 67^{\circ}$};

\end{tikzpicture}
\caption[Pendulum period as a function of amplitude]{The pendulum
period $T(\theta_{0})$ normalised by the small-angle period
$T_{0}=2\pi\sqrt{\ell/g}$, plotted against the amplitude
$\theta_{0}$ in degrees. The small-angle prediction (blue dashed)
is flat at $T/T_{0}=1$. The first amplitude correction
$T/T_{0}\approx 1+\theta_{0}^{2}/16$ (amber dash-dot) tracks the
exact value to within $1\,\%$ up to about $\theta_{0}=70^{\circ}$
and diverges thereafter. The exact value, obtained from the
complete elliptic integral of the first kind $K$, follows the red
curve. The amplitude at which the exact period exceeds the
small-angle period by $5\,\%$ is approximately $67^{\circ}$;
beyond $\theta_{0}\approx 90^{\circ}$ the small-angle formula
under-predicts the period by more than $18\,\%$. The chapter's
project quantifies the crossings; the figure provides the
reference picture against which numerical extractions can be
sanity-checked. Exact values from Abramowitz \& Stegun, table 17.1.}
\label{fig:vol02:ch07:pendulum-period}
\end{figure}


\subsection{Predator-prey and conserved quantities}

The Lotka-Volterra system $\dot{x} = a x - b x y$, $\dot{y} = -c y
+ d x y$ has the conserved quantity
\[
H(x, y) = d x - c \ln x + b y - a \ln y.
\]
The reader can check by computing $\dot{H} = (d - c/x)\dot{x} +
(b - a/y)\dot{y}$ and substituting; the result is identically
zero. Trajectories trace level curves of $H$ in the phase plane,
which are closed orbits.

A naive RK4 integration with $h = 0.01$ over $1{,}000$ time units
shows $H$ drifting by a few percent: the modified equation is not
exactly Hamiltonian. A symplectic integrator preserves $H$ to
machine precision over the same span. The exercise asks the reader
to confirm the difference numerically.

What fails: the textbook predator-prey system is structurally
unstable as a closed-orbit system. Adding a small density-dependent
mortality $-\epsilon x^{2}$ to $\dot{x}$ (logistic correction)
turns the centre into a stable spiral, and the closed orbits decay
to the coexistence equilibrium. Real ecology routinely involves
exactly this correction, and ecologists rarely interpret real
predator-prey time series as the textbook centre.

\subsection{Drug pharmacokinetics: the one-compartment model}

A drug administered intravenously distributes into a single
well-mixed compartment of volume $V$ and is cleared at a rate
proportional to its concentration. Mass balance:
$\dot{m} = -k m$, where $m$ is the drug mass and $k$ has units of
inverse time. The concentration is $C = m/V$. The solution is
$m(t) = m_{0} e^{-k t}$, with $m_{0}$ the dose. The
\engterm{half-life} is $t_{1/2} = \ln 2 / k$.

For acetaminophen (paracetamol) at therapeutic doses, the working
elimination half-life is approximately $2$ to $3$ hours
(pharmacokinetic; the exact value depends on dose, route,
metabolic state, and patient population; clinical references
should be consulted for binding numbers). After five half-lives,
about $97\,\%$ of the dose has been cleared; this is the working
basis of the four-to-six-hour interval between doses. A multi-dose
regimen produces a sequence of decay curves added with the proper
delay; the steady-state peak and trough concentrations follow
geometric-series sums that are exact under the linear assumption.

What fails: the one-compartment linear model assumes
first-order kinetics. Many drugs follow Michaelis-Menten
saturation kinetics at high concentrations, where the elimination
rate plateaus at a maximum velocity instead of growing linearly
with concentration. Ethanol is the canonical example: the
elimination rate is roughly constant rather than proportional to
concentration, so the time to clear a dose grows linearly with the
dose rather than logarithmically.

\subsection{A short history of ODE solution methods}\pathtag{mastery}

The working catalogue of ODE methods is the accumulated answer to
a two-century engineering question: given a differential equation
the integrator could not write in closed form, what step rule
gives a reliable approximate solution? The catalogue's milestones
trace the shifting boundary between what is solvable analytically
and what must be integrated numerically.

\paragraph{Euler 1768.} Leonhard Euler's \emph{Institutiones
Calculi Integralis}, volume I (1768), introduced the polygonal
method that bears his name: from $(t_{n}, y_{n})$, advance to
$(t_{n} + h, y_{n} + h f(t_{n}, y_{n}))$ \cite{hist:v2c07-euler-1768}.
Euler used the method to integrate trajectories for which the
calculus of variations he and Lagrange had developed could not
produce a closed form. The first-order accuracy and the linear
growth of error with the interval length were understood and
accepted as the cost of numerical access to otherwise inaccessible
solutions.

\paragraph{Adams 1855; Bashforth 1883.} John Couch Adams developed
multistep methods in the context of computing lunar and planetary
orbits, formalising the family of explicit predictors that bear
his name and Francis Bashforth's; the methods were published in
Bashforth's 1883 monograph on capillary action, where Adams's
methods were used to integrate the capillary surface ODE
\cite{hist:v2c07-bashforth-adams-1883}. The advance: order of accuracy
increases without proportional growth in right-hand-side
evaluations, by reusing values from previous steps.

\paragraph{Runge 1895; Heun 1900; Kutta 1901.} Carl Runge
introduced the multi-stage methods later named for him in his
1895 paper \cite{hist:v2c07-runge-1895}; Karl Heun and Wilhelm Kutta
extended the family in 1900 and 1901, with Kutta's 1901 paper
introducing the classical fourth-order method that remains the
default reach for non-stiff smooth problems
\cite{hist:v2c07-kutta-1901}. The Runge-Kutta family solved the
multistep methods' starting-value problem (a multistep method
needs initial values from a single-step bootstrap) and supported
adaptive step sizes naturally.

\paragraph{Dahlquist 1956 and the A-stability concept.} Germund
Dahlquist's 1956 paper introduced the theory of $A$-stability and
the systematic distinction between accuracy and stability that
governs the choice of method for stiff problems
\cite{hist:v2c07-dahlquist-1956}. The Dahlquist equivalence theorem
(consistency plus zero-stability implies convergence) and the
later Dahlquist barriers (no explicit linear multistep method of
order $> 2$ is $A$-stable; no $A$-stable multistep method has
order $> 2$) shaped the development of implicit methods.

\paragraph{Gear 1971 and the BDF family.} C.\ William Gear's 1971
monograph \cite{hist:v2c07-gear-1971} introduced the backward
differentiation formulae (BDF) as the canonical implicit family
for stiff problems and made them practical through the DIFSUB
package, the ancestor of LSODE, VODE, and the modern \texttt{ode15s}
and SUNDIALS solvers. Gear's empirical observation that real
engineering problems are routinely stiff, and that the integration
cost is dominated by the implicit linear solves, reshaped the
working scientific-computing stack.

\paragraph{Lorenz 1963.} Edward Lorenz's 1963 paper
\cite{hist:v2c07-lorenz-1963}, integrating a three-equation atmospheric
convection model with a low-order method on an LGP-30 computer at
MIT, produced trajectories whose sensitive dependence on initial
conditions could not be reconciled with the smooth, periodic
behaviour the climatologists had expected. Lorenz's discovery of
deterministic chaos in a working ODE integration recast the
boundary between predictability and unpredictability: not every
ODE that can be integrated produces a forecastable trajectory,
and the working horizon over which a numerical solution carries
information is limited not just by accumulated arithmetic error
but by the system's own Lyapunov exponent. The Lorenz attractor
became the standard teaching example for sensitive dependence and
strange attractors; the methods used to integrate it have improved,
but the qualitative lesson stands.

\paragraph{Modern composite methods.} The current working stack
combines several of the above: adaptive Runge-Kutta methods of
embedded order (Cash-Karp, Dormand-Prince) for non-stiff
problems; BDF and Rosenbrock methods for stiff problems;
geometric (symplectic, energy-preserving) methods for long-time
Hamiltonian integration. Adaptive step controllers detect
stiffness at runtime and switch families automatically. The
working engineer chooses by problem class, not by method
sympathy.

\section{Project}\pathtag{core}

\begin{project}[The pendulum at large amplitudes]
The reader simulates the simple pendulum $\ddot{\theta} +
(g/\ell)\sin\theta = 0$ over the amplitude range
$\theta_{0} = 5^{\circ}$ to $90^{\circ}$ and produces a quantitative
comparison of three predicted periods: the small-angle period
$T_{0} = 2\pi\sqrt{\ell/g}$, the first-correction period $T_{1} =
T_{0}(1 + \theta_{0}^{2}/16)$, and the full numerical period
$T_{\text{num}}$. Per Q55-vol02-ch07 the project track is
simulation with analysis at hazard class none; the reader does not
build a physical pendulum.

\textbf{Hazard class}: none. The project is fully numerical. No
physical apparatus, no biological subject, no mains power. The
reader uses paper and pencil for the analytical work and a
programming environment of choice (Python with SciPy, MATLAB,
Julia, or any environment with an ODE solver) for the simulation.

\textbf{Tools}. Paper, pencil, calculator. A programming
environment with an explicit ODE integrator (RK4 is sufficient;
the working environment's default adaptive solver, e.g.\ SciPy
\path{solve_ivp} with method \path{RK45} or \path{DOP853}, is
preferred at higher amplitudes). A plotting tool. A reference
value for $g$ at the reader's location (the local Earth gravity
varies from about $9.78\,\si{\meter\per\second\squared}$ at the
equator to $9.83\,\si{\meter\per\second\squared}$ at the poles;
the reader chooses a single working value, typically
$9.81\,\si{\meter\per\second\squared}$, and documents the choice).

\textbf{Procedure}.

\begin{enumerate}[leftmargin=*]
  \item \emph{Set up the system}. Choose a length $\ell$, for
    example $\ell = 1\,\si{\meter}$, and the local $g$. Compute
    the small-angle period $T_{0} = 2\pi\sqrt{\ell/g}$. For
    $\ell = 1\,\si{\meter}$, $g = 9.81\,\si{\meter\per\second
    \squared}$: $T_{0} \approx 2.006\,\si{\second}$.

  \item \emph{Convert to a first-order system}. Let $\mathbf{y} =
    (\theta, \dot\theta)^{T}$. Then $\dot{\mathbf{y}} =
    (\dot\theta, -(g/\ell)\sin\theta)^{T}$. Code the right-hand
    side as a function of $t$ and $\mathbf{y}$.

  \item \emph{Numerical period extraction}. For each amplitude
    $\theta_{0}$ in the test grid (e.g.\ $5^{\circ}, 10^{\circ},
    20^{\circ}, 30^{\circ}, 45^{\circ}, 60^{\circ}, 75^{\circ},
    90^{\circ}$), integrate the system from $\theta(0) =
    \theta_{0}$, $\dot\theta(0) = 0$ over enough time to capture
    several periods. Extract the period from the time between
    successive zero crossings of $\dot\theta$ in the negative-to-
    positive direction (or any other consistent crossing
    convention). Use linear interpolation between the two samples
    that bracket the crossing for sub-step resolution. Report the
    average period over at least five oscillations to suppress
    integration drift.

  \item \emph{Compare the three predictions}. Tabulate $T_{0}$
    (small-angle), $T_{1} = T_{0}(1 + \theta_{0}^{2}/16)$ (first
    correction), and $T_{\text{num}}$ (numerical). Compute the
    fractional errors $(T_{0} - T_{\text{num}})/T_{\text{num}}$
    and $(T_{1} - T_{\text{num}})/T_{\text{num}}$ at each
    amplitude.

  \item \emph{Plot}. Produce one figure with three curves: the
    small-angle prediction, the first-correction prediction, and
    the numerical period, all as functions of amplitude
    $\theta_{0}$ from $5^{\circ}$ to $90^{\circ}$. Produce a
    second figure with two curves: the fractional error of each
    approximation against amplitude, on a logarithmic vertical
    axis. Identify the amplitude at which each approximation
    crosses the $1\,\%$ and $5\,\%$ error thresholds.

  \item \emph{Step-size convergence study}. At
    $\theta_{0} = 60^{\circ}$, repeat the numerical period
    extraction at four step sizes $h$ in geometric progression
    (e.g.\ $h = 0.01, 0.005, 0.0025, 0.00125\,\si{\second}$).
    Plot the period error against $h$ on log-log axes. Confirm
    that the error scales as $h^{4}$ for RK4 (slope $4$ on the
    log-log plot) and as $h$ for forward Euler (slope $1$, if
    Euler is included for comparison).

  \item \emph{Conservation diagnostic}. The pendulum's mechanical
    energy $E = \tfrac{1}{2}\ell^{2}\dot\theta^{2} + g\ell(1 -
    \cos\theta)$ (per unit mass) is conserved by the exact
    dynamics. For each amplitude, compute the maximum drift in
    $E$ over the simulation interval. Report the drift as a
    fraction of the initial $E$. Compare RK4 against a symplectic
    integrator (leapfrog or velocity-Verlet) at the same step
    size.
\end{enumerate}

\textbf{Record}. The deliverable is a notebook (Jupyter, MATLAB
live script, or Pluto) or a script-plus-report containing: the
right-hand-side definition; the period-extraction routine; the
period-versus-amplitude tabulation and plot; the error-versus-
amplitude plot; the step-size convergence study and plot; the
conservation diagnostic and plot; and a 1500-word reflection
addressing the questions below.

\textbf{Reflection (1500 words)}. The reader writes a 1500-word
narrative addressing:

\begin{itemize}[leftmargin=*]
  \item At what amplitude does the small-angle period exceed the
    numerical period by more than $1\,\%$? By more than $5\,\%$?
    Compare the answers to the first-correction prediction's
    crossing of the same thresholds, and explain the gap in
    physical terms.
  \item How does the fourth-order convergence of RK4 manifest in
    the period-error plot, and what limits the achievable
    precision at the smallest step size? (Hint: floating-point
    roundoff begins to dominate truncation error at sufficiently
    small $h$; the reader should observe this floor.)
  \item How does the energy drift of RK4 compare with the energy
    drift of the symplectic integrator at the same step size?
    What does the comparison say about the choice of integrator
    for problems where the long-term qualitative behaviour
    matters?
  \item One physical effect (air resistance, finite rod mass,
    pivot friction, bob deformation) that the simple-pendulum
    model neglects but a real pendulum experiences. State, in
    physical terms, the order of magnitude of the effect for a
    one-metre laboratory pendulum, and identify whether the
    correction is larger or smaller than the leading
    nonlinearity correction at $\theta_{0} = 30^{\circ}$.
  \item What numerical and analytical methods buy each other.
    The numerical method gives the answer the analytical method
    cannot, and the analytical method gives the structure
    (small-angle period as the leading order, the
    $\theta_{0}^{2}/16$ correction as the next term, the
    elliptic-integral exact answer as the closed form) that the
    numerical method by itself does not surface.
\end{itemize}

\textbf{Time}. Setup and small-angle implementation: one to two
hours. Period extraction across the amplitude grid: one to two
hours. Convergence study and conservation diagnostic: one to two
hours. Reflection: two to three hours. Total reader time
approximately five to nine hours.

\textbf{Phases}. The reader works through the project in five
explicit phases. Phase 1 (analytical preparation, $\sim 1$ hour):
derive the linearisation, the first amplitude correction, and the
energy expression by hand. Phase 2 (implementation, $\sim 1$
hour): code the right-hand side and the RK4 step; verify against
the closed-form exponential decay of a damped linear oscillator
before tackling the pendulum. Phase 3 (period extraction, $\sim
2$ hours): run the amplitude sweep and the zero-crossing
extractor; tabulate the three predictions. Phase 4 (convergence
and conservation, $\sim 2$ hours): repeat the period extraction at
multiple step sizes; produce the log-log convergence plot; compute
the energy drift for RK4 and a symplectic integrator. Phase 5
(reflection, $\sim 2$ hours): write the $1500$-word narrative.

\textbf{Success criteria}. The deliverable meets the standard when
(i) the small-angle prediction matches $T_{\text{num}}$ to within
$0.1\,\%$ at $\theta_{0} = 5^{\circ}$ and disagrees by more than
$5\,\%$ at $\theta_{0} = 90^{\circ}$, with the disagreement growing
monotonically with amplitude; (ii) the first-correction prediction
$T_{1} = T_{0}(1 + \theta_{0}^{2}/16)$ agrees with $T_{\text{num}}$
to within $1\,\%$ up to about $\theta_{0} = 70^{\circ}$, beyond
which higher-order corrections kick in (see \cref{fig:vol02:ch07:pendulum-period});
(iii) the log-log convergence plot for RK4 has slope close to $4$
in the truncation-dominated regime and flattens at small $h$ where
floating-point roundoff begins to dominate; (iv) the energy drift
of the symplectic integrator is smaller than the energy drift of
RK4 by at least one order of magnitude at the same step size
over a multi-period run.

\textbf{Common failure modes}. Three failures recur and the reader
should know them in advance. First, the small-angle period
\emph{matches} the numerical period across the whole amplitude
range, despite the analysis predicting growing disagreement. The
cause is usually a bug in the right-hand side (e.g.\ using $\theta$
instead of $\sin\theta$, which silently keeps the system linear).
The diagnostic is to plot the right-hand side at $\theta_{0} =
60^{\circ}$ and compare against a hand calculation. Second, the
energy drift of RK4 \emph{decreases} over time, in apparent
violation of the modified-equation prediction. The cause is
usually a sign error in the right-hand side that makes the
modified equation Hamiltonian by accident; the diagnostic is to
reverse the time integration and check that the trajectory retraces
the forward run. Third, the convergence plot for RK4 shows slope
close to $1$ rather than $4$. The cause is usually a step-size
adaptation that has overridden the requested fixed step, or a
right-hand-side function that uses the wrong arguments at the
intermediate Runge--Kutta stages. The diagnostic is to print the
actual step size taken by the integrator at the end of the run.
\end{project}

\section{Exercises}

The exercises run from first-order quadrature through forced
response, systems, phase-plane diagnosis, numerical solver tuning,
and the failure-analysis case studies that close the chapter.

\subsection*{Calculation}\pathtag{core}

\begin{exercise}
Solve the IVP $\dot{y} + 2 y = e^{-t}$, $y(0) = 1$, by the
integrating-factor method. Verify by substitution.

\paragraph{Solution sketch.} Standard form: $p(t) = 2$, $q(t) =
e^{-t}$. Integrating factor $\mu(t) = e^{2 t}$. Multiplying the
ODE by $\mu$ gives $\frac{d}{dt}(e^{2 t} y) = e^{2 t} \cdot e^{-t}
= e^{t}$. Integrating from $0$ to $t$: $e^{2 t} y(t) - 1 = e^{t}
- 1$, hence $e^{2 t} y(t) = e^{t}$, so $y(t) = e^{-t}$. Verify:
$\dot{y} = -e^{-t}$; $\dot{y} + 2 y = -e^{-t} + 2 e^{-t} = e^{-t}$,
matching the right-hand side, and $y(0) = 1$ matches the initial
condition.
\end{exercise}

\begin{exercise}
Solve the separable IVP $\dot{y} = y \cos t$, $y(0) = 2$. Identify
the period and amplitude of the oscillating solution envelope.

\paragraph{Solution sketch.} Separate: $dy/y = \cos t \, dt$.
Integrate: $\ln |y| = \sin t + C$, hence $y(t) = A e^{\sin t}$.
The initial condition $y(0) = 2$ gives $A = 2$, so $y(t) = 2
e^{\sin t}$. The solution is not oscillating in the linear sense;
it modulates between extremes. The maximum value $2 e^{1}
\approx 5.44$ occurs at $t = \pi/2 + 2\pi k$; the minimum
$2 e^{-1} \approx 0.74$ at $t = 3\pi/2 + 2\pi k$. The period of
the modulation is $2\pi$. The envelope is bounded; the solution
returns to its starting value $2$ at $t = 0, \pi, 2\pi, \ldots$
(whenever $\sin t = 0$). The reader who fits this form to data
should recognise that the multiplicative-noise drive produces a
log-normal envelope, not a sinusoidal one.
\end{exercise}

\begin{exercise}
For the second-order ODE $\ddot{x} + 4 \dot{x} + 13 x = 0$ with
$x(0) = 1$, $\dot{x}(0) = 0$, compute the natural frequency
$\omega_{0}$, the damping ratio $\zeta$, the damped frequency
$\omega_{d}$, and write the explicit solution. Identify the
half-life of the amplitude envelope.

\paragraph{Solution sketch.} Identify $m = 1$, $c = 4$, $k = 13$.
Natural frequency $\omega_{0} = \sqrt{k/m} = \sqrt{13} \approx
3.606\,\si{\radian\per\second}$. Damping ratio $\zeta = c/(2
\sqrt{m k}) = 4/(2 \sqrt{13}) = 2/\sqrt{13} \approx 0.555$.
Discriminant $c^{2} - 4 m k = 16 - 52 = -36 < 0$: underdamped.
Damped frequency $\omega_{d} = \omega_{0}\sqrt{1 - \zeta^{2}} =
\sqrt{13} \cdot \sqrt{1 - 4/13} = \sqrt{13 - 4} = 3\,\si{\radian
\per\second}$. Decay rate $\zeta \omega_{0} = 2$. The general
solution is $x(t) = e^{-2 t}(A \cos(3 t) + B \sin(3 t))$. From
$x(0) = 1$: $A = 1$. From $\dot{x}(0) = 0$: $-2 A + 3 B = 0$,
hence $B = 2/3$. Therefore $x(t) = e^{-2 t}(\cos 3 t +
(2/3) \sin 3 t)$. The amplitude envelope $\sqrt{A^{2} + B^{2}}\,
e^{-\zeta \omega_{0} t} = \sqrt{1 + 4/9}\, e^{-2 t} \approx
1.202\, e^{-2 t}$ has half-life $t_{1/2} = \ln 2 / (\zeta
\omega_{0}) = \ln 2 / 2 \approx 0.347\,\si{\second}$.
\end{exercise}

\begin{exercise}
For an RC circuit with $R = 1\,\si{\kilo\ohm}$ and $C =
1\,\si{\micro\farad}$ initially at $10\,\si{\volt}$, compute the
voltage at $t = 1\,\si{\milli\second}$ and the time required for
the voltage to fall below $1\,\si{\milli\volt}$.

\paragraph{Solution sketch.} Time constant $\tau = RC = 10^{3}
\times 10^{-6} = 10^{-3}\,\si{\second} = 1\,\si{\milli\second}$.
The discharge solution is $v(t) = 10\, e^{-t/\tau}\,\si{\volt}$.
At $t = 1\,\si{\milli\second}$: $v = 10/e \approx 3.68\,\si{
\volt}$. For $v < 1\,\si{\milli\volt} = 10^{-3}\,\si{\volt}$,
require $e^{-t/\tau} < 10^{-4}$, i.e.\ $t/\tau > \ln 10^{4} = 4
\ln 10 \approx 9.21$, so $t > 9.21\,\si{\milli\second}$. The
working memory: each time constant cuts the voltage by a factor
of $e \approx 2.72$; four orders of magnitude take about $9.2$
time constants.
\end{exercise}

\begin{exercise}
A drug administered as a single $500\,\si{\milli\gram}$ IV bolus
into a one-compartment model has elimination rate constant
$k = 0.231\,\si{\per\hour}$. Compute the concentration as a
fraction of initial concentration at $t = 1, 3, 6,
12\,\si{\hour}$. Compute the half-life. Estimate the time at
which the concentration falls below $5\,\%$ of its initial value.

\paragraph{Solution sketch.} The solution is $C(t)/C(0) = e^{-k
t}$. Tabulate: $e^{-0.231} \approx 0.794$; $e^{-0.693} \approx
0.500$; $e^{-1.386} \approx 0.250$; $e^{-2.772} \approx 0.0625$.
The half-life is $t_{1/2} = \ln 2 / k = 0.693 / 0.231 \approx
3.00\,\si{\hour}$; the $k$ was chosen to give a clean
three-hour half-life. For $C/C(0) < 0.05$, require $k t > \ln 20
\approx 3.00$, hence $t > 13.0\,\si{\hour}$. After about four
half-lives the drug is essentially cleared from the single
compartment under the linear-first-order model.
\end{exercise}

\begin{exercise}
Compute the matrix exponential $e^{A t}$ for $A = \begin{pmatrix}
-1 & 1 \\ 0 & -2 \end{pmatrix}$. Verify by checking that
$\frac{d}{dt}e^{A t}\bigl|_{t=0} = A$.

\paragraph{Solution sketch.} The matrix is upper triangular with
distinct eigenvalues $\lambda_{1} = -1$, $\lambda_{2} = -2$.
Diagonalise: $A = P D P^{-1}$ with $D = \operatorname{diag}(-1,
-2)$. Eigenvectors: for $\lambda_{1} = -1$, $(A - (-1)I)v = 0$
gives $v_{1} = (1, 0)^{T}$. For $\lambda_{2} = -2$, $(A + 2 I) v
= 0$ gives $v_{2} = (1, -1)^{T}$ (any non-zero scaling works).
So $P = \begin{pmatrix} 1 & 1 \\ 0 & -1 \end{pmatrix}$ and $P^{-1}
= \begin{pmatrix} 1 & 1 \\ 0 & -1 \end{pmatrix}$ (the inverse
equals $P$ here by direct check). Then
\[
e^{A t} = P \begin{pmatrix} e^{-t} & 0 \\ 0 & e^{-2 t}
\end{pmatrix} P^{-1} = \begin{pmatrix} e^{-t} & e^{-t} - e^{-2 t}
\\ 0 & e^{-2 t} \end{pmatrix}.
\]
Verify: $\left.\frac{d}{dt}e^{A t}\right|_{t = 0} =
\begin{pmatrix} -1 & 1 \\ 0 & -2 \end{pmatrix} = A$. \checkmark
\end{exercise}

\begin{exercise}
For the predator-prey system with $a = 1$, $b = 0.5$, $c = 0.75$,
$d = 0.25$, locate the coexistence equilibrium and compute the
linearisation's eigenvalues. State the predicted oscillation
period from the imaginary part of the eigenvalues.

\paragraph{Solution sketch.} Coexistence equilibrium $(x^{*},
y^{*}) = (c/d, a/b) = (3, 2)$. At this point, the Jacobian is
\[
J = \begin{pmatrix} a - b y & -b x \\ d y & -c + d x
\end{pmatrix}\bigg|_{(3,2)} = \begin{pmatrix} 0 & -1.5 \\ 0.5 & 0
\end{pmatrix}.
\]
Trace $T = 0$, determinant $D = 0.75 > 0$, discriminant $T^{2} -
4 D = -3 < 0$: pure imaginary eigenvalues, centre (in the
linearisation). Eigenvalues $\lambda = \pm i \sqrt{D} = \pm i
\sqrt{0.75} \approx \pm 0.866 i$. Predicted oscillation period $T
= 2\pi / |\operatorname{Im} \lambda| = 2\pi/\sqrt{0.75} \approx
7.26$ time units. The exact period of the full nonlinear closed
orbit depends on amplitude; the linearised prediction is valid
for orbits close to the equilibrium and grows by a few percent
for amplitudes comparable to the equilibrium values.
\end{exercise}

\begin{exercise}
For $\dot{y} = -10 y + \cos t$ with $y(0) = 0$, write down the
steady-state response amplitude and phase using the
forced-oscillator formula adapted to first order. Verify that the
homogeneous transient decays with time constant $0.1\,\si
{\second}$.

\paragraph{Solution sketch.} Standard form: $\dot{y} + 10 y =
\cos t$. Try $y_{p}(t) = \operatorname{Re}(\hat{Y} e^{i t})$ with
$\hat{Y}$ complex. Substituting and matching: $i \hat{Y} + 10
\hat{Y} = 1$, so $\hat{Y} = 1/(10 + i)$. Magnitude $|\hat{Y}| =
1/\sqrt{101} \approx 0.0995$. Argument $\arg \hat{Y} = -\arctan
(1/10) \approx -0.0997\,\si{\radian} \approx -5.71^{\circ}$.
Steady-state response: $y_{p}(t) \approx 0.0995 \cos(t -
0.0997)$. The homogeneous mode is $e^{-10 t}$, with time
constant $\tau = 1/10 = 0.1\,\si{\second}$. \checkmark The full
solution is $y(t) = y_{h}(t) + y_{p}(t)$ with $y_{h}(0) = -y_{p}(0)$
to enforce $y(0) = 0$; the transient decays in a few tenths of
a second, and the steady-state response is the long-time behaviour.
\end{exercise}

\begin{exercise}
A mass-spring-dashpot system with $m = 2\,\si{\kilogram}$, $c =
3\,\si{\newton\second\per\meter}$, $k = 50\,\si{\newton\per\meter}$
is driven by $F(t) = 5\cos(\omega t)\,\si{\newton}$. Compute the
steady-state amplitude as a function of $\omega$, identify the
resonant frequency, and compute the amplitude at the resonant
frequency. Compute the quality factor $Q$.

\paragraph{Solution sketch.} Natural frequency $\omega_{0} =
\sqrt{k/m} = \sqrt{25} = 5\,\si{\radian\per\second}$. Damping
ratio $\zeta = c/(2 \sqrt{m k}) = 3/(2 \sqrt{100}) = 0.15$.
Quality factor $Q = 1/(2 \zeta) = 3.33$. Steady-state amplitude
\[
A(\omega) = \frac{F_{0}}{\sqrt{(k - m \omega^{2})^{2} + (c
\omega)^{2}}} = \frac{5}{\sqrt{(50 - 2 \omega^{2})^{2} + 9
\omega^{2}}}.
\]
Resonant frequency (peak of $A$) is $\omega_{r} = \omega_{0}
\sqrt{1 - 2 \zeta^{2}} = 5 \sqrt{1 - 0.045} \approx 4.886\,\si{
\radian\per\second}$. At resonance, $A(\omega_{r}) \approx
F_{0}/(c \omega_{0}) \times 1/\sqrt{1 - \zeta^{2}} = 5/(3 \times
5) / \sqrt{1 - 0.0225} \approx 0.337\,\si{\meter}$. Compare to
the static deflection $F_{0}/k = 5/50 = 0.10\,\si{\meter}$: the
resonance amplifies by a factor of about $3.4$, close to $Q$ as
expected for moderate damping.
\end{exercise}

\begin{exercise}
For the system $\dot{x} = -x + 2 y$, $\dot{y} = -3 x - 4 y$,
write the matrix $A$, compute its eigenvalues and eigenvectors,
and write the general solution. Identify whether the origin is a
stable node, a stable spiral, a saddle, or otherwise.

\paragraph{Solution sketch.} $A = \begin{pmatrix} -1 & 2 \\ -3 &
-4 \end{pmatrix}$. Trace $T = -5$, determinant $D = (-1)(-4) -
(2)(-3) = 4 + 6 = 10$. Discriminant $T^{2} - 4 D = 25 - 40 = -15
< 0$: complex eigenvalues; $T < 0$, $D > 0$: stable spiral.
Eigenvalues $\lambda = (T \pm i \sqrt{-T^{2} + 4 D})/2 = (-5 \pm
i \sqrt{15})/2$. Real part $-5/2 = -2.5$ (decay rate); imaginary
part $\sqrt{15}/2 \approx 1.936$ (oscillation frequency).
Eigenvector for $\lambda_{+}$: solve $(A - \lambda_{+} I) v = 0$.
First-row equation: $(-1 - \lambda_{+}) v_{1} + 2 v_{2} = 0$,
giving $v_{2}/v_{1} = (1 + \lambda_{+})/2 = (-3 + i\sqrt{15})/4$.
General real solution: $\mathbf{y}(t) = e^{-2.5 t}(\mathbf{c}_{1}
\cos(\sqrt{15}\, t/2) + \mathbf{c}_{2} \sin(\sqrt{15}\, t/2))$,
with real vector constants $\mathbf{c}_{i}$ pinned by the initial
condition. Trajectories spiral into the origin counter-clockwise.
\end{exercise}

\subsection*{Derivation}\pathtag{standard}

\begin{exercise}
Derive the integrating-factor formula $\mu(t) = \exp(\int p(t)\,dt)$
for the standard first-order linear ODE $\dot{y} + p(t) y = q(t)$.
Show, by direct substitution, that any antiderivative for the
integral defines an integrating factor, with the constant of
integration cancelling in the final solution.

\paragraph{Solution sketch.} Multiply the ODE by an unknown
$\mu(t)$: $\mu \dot{y} + \mu p y = \mu q$. We want the left side
to equal $\frac{d}{dt}(\mu y) = \dot{\mu} y + \mu \dot{y}$, so
require $\dot{\mu} = \mu p$, i.e.\ $d \mu/\mu = p \, dt$. Integrate:
$\ln \mu = \int p \, dt + C_{0}$, giving $\mu = e^{C_{0}}
\exp(\int p \, dt) = A \exp(\int p \, dt)$ for arbitrary positive
$A$. The general solution becomes $\mu y = \int \mu q \, dt + C_{1}$,
i.e.\ $A \exp(\int p) \, y = \int A \exp(\int p) \, q \, dt + C_{1}$,
which simplifies to $y = \exp(-\int p)[\int \exp(\int p) \, q \,
dt + C_{1}/A]$. The factor $A$ multiplies both sides of the
integral equation and divides out; only the constant $C_{1}/A$
appears, and it is set by the initial condition. The antiderivative
choice (the constant $C_{0}$) thus does not affect the solution.
\end{exercise}

\begin{exercise}
Derive the characteristic polynomial $m \lambda^{2} + c \lambda +
k = 0$ for the second-order linear constant-coefficient ODE by
substituting $x = e^{\lambda t}$ into $m \ddot{x} + c \dot{x} +
k x = 0$. Show that the discriminant $c^{2} - 4 m k$ classifies
the homogeneous behaviour into the three regimes (overdamped,
critically damped, underdamped) and identify the boundary case.

\paragraph{Solution sketch.} Substitute $x = e^{\lambda t}$:
$\dot{x} = \lambda e^{\lambda t}$, $\ddot{x} = \lambda^{2}
e^{\lambda t}$. The ODE becomes $e^{\lambda t}(m \lambda^{2} + c
\lambda + k) = 0$. Since $e^{\lambda t} \neq 0$, we require
$m \lambda^{2} + c \lambda + k = 0$, the characteristic polynomial.
The roots are $\lambda_{\pm} = (-c \pm \sqrt{c^{2} - 4 m k})/(2 m)$.
For $c^{2} > 4 m k$ (discriminant positive): two distinct real
negative roots; overdamped. For $c^{2} = 4 m k$ (discriminant
zero): one repeated real root $\lambda = -c/(2 m)$, with the
second linearly-independent solution $t e^{\lambda t}$ arising
from a Jordan block; critically damped. For $c^{2} < 4 m k$
(discriminant negative): complex conjugate roots with negative
real part; underdamped, with oscillation. The boundary case is
$c^{2} = 4 m k$, i.e.\ $c = 2 \sqrt{m k}$, which sets $\zeta = 1$.
\end{exercise}

\begin{exercise}
Derive the steady-state amplitude and phase for the forced linear
oscillator $m \ddot{x} + c \dot{x} + k x = F_{0}\cos(\omega t)$
using complex exponentials. Show that the amplitude is maximised
at $\omega_{r} = \omega_{0}\sqrt{1 - 2\zeta^{2}}$ and that the
peak amplitude is approximately $F_{0}/(c\omega_{0})$ for light
damping.

\paragraph{Solution sketch.} Use complex exponentials: $F(t) =
\operatorname{Re}(F_{0} e^{i \omega t})$; try $x_{p}(t) =
\operatorname{Re}(\hat{X} e^{i \omega t})$. Substitute: $(-m
\omega^{2} + i c \omega + k) \hat{X} = F_{0}$, hence
\[
\hat{X} = \frac{F_{0}}{(k - m \omega^{2}) + i c \omega},
\quad |\hat{X}| = \frac{F_{0}}{\sqrt{(k - m \omega^{2})^{2} + (c
\omega)^{2}}}.
\]
To maximise $|\hat{X}|$ as a function of $\omega$, minimise the
denominator's square: $D(\omega) = (k - m \omega^{2})^{2} + (c
\omega)^{2}$. Set $dD/d\omega = 0$: $-4 m (k - m \omega^{2})
\omega + 2 c^{2} \omega = 0$. For $\omega \neq 0$, $2 m (k - m
\omega^{2}) = c^{2}$, so $\omega^{2} = (k - c^{2}/(2 m))/m =
\omega_{0}^{2} - c^{2}/(2 m^{2})$. Using $\zeta = c/(2 \sqrt{m k})$,
$c^{2}/(2 m^{2}) = 2 \zeta^{2} \omega_{0}^{2}$, hence $\omega_{r}
= \omega_{0} \sqrt{1 - 2 \zeta^{2}}$. At $\omega = \omega_{0}$,
$|\hat{X}|_{\omega = \omega_{0}} = F_{0}/(c \omega_{0})$, which
approximates the peak to leading order in $\zeta$ for light
damping.
\end{exercise}

\begin{exercise}
Show by direct calculation that $H(x, y) = d x - c \ln x + b y -
a \ln y$ is conserved along the flow of the Lotka-Volterra system
$\dot{x} = a x - b x y$, $\dot{y} = -c y + d x y$. Sketch the
level curves in the $(x, y)$ plane for one choice of parameters.

\paragraph{Solution sketch.} Compute the time derivative of $H$
along trajectories:
\[
\dot{H} = \left(d - \frac{c}{x}\right) \dot{x} + \left(b - \frac{a}
{y}\right) \dot{y}.
\]
Substitute $\dot{x} = x(a - b y)$ and $\dot{y} = y(d x - c)$:
\[
\dot{H} = \left(d - \frac{c}{x}\right) x (a - b y) + \left(b -
\frac{a}{y}\right) y (d x - c).
\]
Expand:
\[
\dot{H} = (d x - c)(a - b y) + (b y - a)(d x - c) = (d x - c)
[(a - b y) + (b y - a)] = 0.
\]
The bracket vanishes identically, so $\dot{H} = 0$ along
trajectories. Level curves of $H$ are closed orbits in the
positive quadrant; the equilibrium $(c/d, a/b)$ is the minimum of
$H$ (a stable centre in the linearisation). For $a = b = c = d =
1$, the minimum is at $(1, 1)$ and orbits surround it.
\end{exercise}

\begin{exercise}
Derive the first amplitude correction to the simple-pendulum
period. Starting from $\ddot\theta + \omega_{0}^{2}\sin\theta =
0$, expand $\sin\theta = \theta - \theta^{3}/6 + O(\theta^{5})$,
assume a solution of the form $\theta(t) =
\theta_{0}\cos(\omega t)$ with $\omega = \omega_{0}(1 +
\alpha\theta_{0}^{2} + O(\theta_{0}^{4}))$, and determine $\alpha$
by requiring the equation to balance at order $\theta_{0}^{3}$.
The result is $\alpha = -1/16$ for the angular frequency, which
gives $T = T_{0}(1 + \theta_{0}^{2}/16 + \cdots)$ for the period.

\paragraph{Solution sketch.} The full derivation uses the
Lindstedt--Poincar\'e perturbation method. The ansatz $\theta =
\theta_{0} \cos(\omega t)$ is a trial-form simplification; the
honest derivation needs $\theta = \theta_{0} \cos(\omega t) +
\theta_{0}^{3} u_{1}(\omega t) + \cdots$ to absorb the harmonics
generated by $\sin \theta$. Substituting into $\ddot\theta +
\omega_{0}^{2} \sin\theta = 0$ with $\sin\theta = \theta - \theta^{3}/6
+ \cdots$ produces at $O(\theta_{0})$ the linear equation
($\omega^{2} = \omega_{0}^{2}$ at leading order) and at $O(\theta_{0}^{3})$
the inhomogeneous equation
\[
u_{1}'' + u_{1} = -2 \alpha \cos\tau + \frac{1}{6} \cos^{3} \tau,
\]
with $\tau = \omega t$ and the right-hand side carrying both a
$\cos\tau$ term (from the $\omega$ correction) and a $\cos^{3}\tau$
term (from the cubic in $\sin$). Using $\cos^{3}\tau =
(3 \cos\tau + \cos 3\tau)/4$ and demanding that the secular
(resonant) $\cos\tau$ coefficient vanish (otherwise $u_{1}$
grows linearly in $\tau$, violating boundedness) gives the
balance $-2 \alpha + (1/6)(3/4) = 0$, hence $\alpha = 1/16$ as a
correction to $\omega^{2}/\omega_{0}^{2}$, i.e.\ $\omega = \omega_{0}
\sqrt{1 - \theta_{0}^{2}/8} \approx \omega_{0}(1 -
\theta_{0}^{2}/16)$. Period $T = 2\pi/\omega \approx T_{0}(1 +
\theta_{0}^{2}/16)$. (The sign of $\alpha$ in $\omega$ is
negative, but the correction to $T$ is positive because $T \propto
1/\omega$.) See \cref{fig:vol02:ch07:pendulum-period} for the
graphical comparison against the exact (elliptic-integral) value.
\end{exercise}

\begin{exercise}
Show that forward Euler applied to $\dot{y} = \lambda y$ produces
$y_{n+1} = (1 + h\lambda) y_{n}$, and that implicit Euler
produces $y_{n+1} = y_{n}/(1 - h\lambda)$. State the stability
region in the complex $h\lambda$ plane for each method.

\paragraph{Solution sketch.} Forward Euler: $y_{n+1} = y_{n} + h
f(t_{n}, y_{n}) = y_{n} + h \lambda y_{n} = (1 + h \lambda) y_{n}$.
Implicit Euler: $y_{n+1} = y_{n} + h f(t_{n+1}, y_{n+1}) = y_{n}
+ h \lambda y_{n+1}$, i.e.\ $(1 - h \lambda) y_{n+1} = y_{n}$,
hence $y_{n+1} = y_{n}/(1 - h \lambda)$. The iterates are
bounded iff the amplification factor has modulus at most $1$.
Forward Euler: $|1 + h \lambda| \leq 1$, the closed disc of
radius $1$ centred at $-1$ in the $z = h \lambda$ plane. Implicit
Euler: $|1/(1 - h \lambda)| \leq 1$, equivalently $|1 - h
\lambda| \geq 1$, the exterior of the closed disc of radius $1$
centred at $+1$. This includes the entire left half-plane,
making implicit Euler \emph{A-stable}; forward Euler is not. See
\cref{fig:vol02:ch07:stiffness} for the visualisation.
\end{exercise}

\begin{exercise}
Derive the local truncation error of forward Euler from the
Taylor expansion $y(t + h) = y(t) + h\dot{y}(t) +
(h^{2}/2)\ddot{y}(t) + O(h^{3})$. Show that the leading error per
step is $(h^{2}/2)\ddot{y}(t)$, and that the global error after
$N = T/h$ steps is $O(h)$ on a finite interval.

\paragraph{Solution sketch.} Forward Euler advances $y_{n+1} =
y_{n} + h f(t_{n}, y_{n})$, while the exact solution satisfies
$y(t_{n} + h) = y(t_{n}) + h \dot{y}(t_{n}) + (h^{2}/2)\ddot{y}
(t_{n}) + O(h^{3})$. The local truncation error (one step,
starting from the exact value) is the difference: $y(t_{n+1}) -
[y(t_{n}) + h \dot{y}(t_{n})] = (h^{2}/2) \ddot{y}(t_{n}) +
O(h^{3})$. On a finite interval $[t_{0}, T]$, the global error
propagates: errors accumulate over $N = (T - t_{0})/h$ steps, and
each step's error is amplified by the local Lipschitz constant of
$f$. The standard Gronwall argument bounds the accumulated error
by $C h$ for some constant $C$ depending on $f$ and the interval
length; the leading-order behaviour is $\|e_{N}\| \leq (e^{L
(T - t_{0})} - 1) / L \cdot (h/2) \max |\ddot{y}|$, where $L$ is
the Lipschitz constant. The global error is $O(h)$: forward Euler
is first-order.
\end{exercise}

\begin{exercise}
Derive the variation-of-parameters formula
$\mathbf{y}(t) = e^{A t}\mathbf{y}_{0} + \int_{0}^{t} e^{A(t -
s)}\mathbf{q}(s)\,ds$ for the driven linear system
$\dot{\mathbf{y}} = A\mathbf{y} + \mathbf{q}(t)$ with
$\mathbf{y}(0) = \mathbf{y}_{0}$. Hint: try $\mathbf{y}(t) = e^{A
t}\mathbf{c}(t)$ and determine $\dot{\mathbf{c}}$.

\paragraph{Solution sketch.} Substitute $\mathbf{y}(t) = e^{A t}
\mathbf{c}(t)$ into $\dot{\mathbf{y}} = A \mathbf{y} + \mathbf{q}$:
\[
A e^{A t} \mathbf{c} + e^{A t} \dot{\mathbf{c}} = A e^{A t}
\mathbf{c} + \mathbf{q}.
\]
The terms involving $A \mathbf{c}$ cancel, giving $e^{A t}
\dot{\mathbf{c}} = \mathbf{q}(t)$, so $\dot{\mathbf{c}}(t) = e^{-A
t} \mathbf{q}(t)$ (using $(e^{A t})^{-1} = e^{-A t}$). Integrate
from $0$ to $t$: $\mathbf{c}(t) = \mathbf{c}(0) + \int_{0}^{t}
e^{-A s} \mathbf{q}(s) \, ds$. The initial condition $\mathbf{y}(0)
= \mathbf{y}_{0}$ gives $\mathbf{c}(0) = \mathbf{y}_{0}$
(since $e^{A \cdot 0} = I$). Substituting back:
\[
\mathbf{y}(t) = e^{A t} \mathbf{y}_{0} + e^{A t} \int_{0}^{t}
e^{-A s} \mathbf{q}(s) \, ds = e^{A t} \mathbf{y}_{0} +
\int_{0}^{t} e^{A (t - s)} \mathbf{q}(s) \, ds,
\]
the variation-of-parameters formula. The factorisation $e^{A t}
e^{-A s} = e^{A (t - s)}$ uses commutativity of $A$ with itself.
\end{exercise}

\subsection*{Estimation}\pathtag{core}

\begin{exercise}
Estimate, before computing, the time for the voltage on a $C =
1\,\si{\micro\farad}$ capacitor discharging through a $R =
1\,\si{\mega\ohm}$ resistor to fall from $10\,\si{\volt}$ to
$0.1\,\si{\volt}$. Compare to the formal answer.

\paragraph{Solution sketch.} Time constant $\tau = RC = 10^{6}
\times 10^{-6} = 1\,\si{\second}$. The ratio $10/0.1 = 100$ takes
$\ln 100 = 2 \ln 10 \approx 4.6$ time constants. Estimate: about
$4.6\,\si{\second}$. Formal: $0.1 = 10 e^{-t/1}$, $t = \ln 100
= 4.605\,\si{\second}$. The estimate matches the formal answer to
all reported precision; the discharge of two orders of magnitude
takes roughly five time constants.
\end{exercise}

\begin{exercise}
Estimate the natural frequency of a simple pendulum of length
$25\,\si{\centi\meter}$. Estimate the period at small amplitude
and at $\theta_{0} = 60^{\circ}$. Compare the two and identify
the fractional change attributable to amplitude.
\paragraph{Solution sketch.} Natural angular frequency $\omega_{0}
= \sqrt{g/\ell} = \sqrt{9.81/0.25} \approx 6.26\,\si{\radian\per
\second}$. Period at small amplitude $T_{0} = 2\pi/\omega_{0}
\approx 1.004\,\si{\second}$. At $\theta_{0} = 60^{\circ} = \pi/3
\approx 1.047\,\si{\radian}$, the first correction gives $T_{1}
= T_{0}(1 + \theta_{0}^{2}/16) = T_{0}(1 + 1.097/16) \approx T_{0}
\times 1.0685$. The exact (elliptic) value from
\cref{fig:vol02:ch07:pendulum-period} is $T/T_{0} \approx 1.073$,
so the first correction is within $0.5\,\%$ of the exact answer.
Numerical period $\approx 1.077\,\si{\second}$. The fractional
increase from small-amplitude is about $7.3\,\%$, large enough
that a $60^{\circ}$ amplitude meaningfully changes a pendulum-
clock's rate: the rate change is the inverse, about $-7.3\,\%$,
$\sim 6000\,\si{\second}$ per day, or one hour fifty minutes per
day.
\end{exercise}

\begin{exercise}
Estimate the quality factor $Q$ of a guitar string of mass per
unit length $\mu = 5 \times 10^{-4}\,\si{\kilogram\per\meter}$,
length $0.65\,\si{\meter}$, tension $80\,\si{\newton}$, with the
sustained sound decaying audibly over $\sim 5\,\si{\second}$.
State the assumption that converts an audible decay time into a
$1/e$ amplitude time. The fundamental frequency follows from $f =
(1/2L)\sqrt{T/\mu}$.

\paragraph{Solution sketch.} Fundamental frequency $f_{0} =
(1/(2 \times 0.65)) \sqrt{80/(5 \times 10^{-4})} = (1/1.30) \times
400 \approx 308\,\si{\hertz}$. Angular frequency $\omega_{0} = 2
\pi f_{0} \approx 1933\,\si{\radian\per\second}$. Assume the
audible decay time of $5\,\si{\second}$ corresponds to amplitude
decay by a factor of about $1/100$ (the loudness drops by
$40\,\si{\decibel}$); the $1/e$ amplitude time is then
$\tau_{1/e} = 5/\ln 100 \approx 1.09\,\si{\second}$. Damping
ratio $\zeta \approx 1/(\omega_{0} \tau_{1/e}) = 1/(1933 \times
1.09) \approx 4.7 \times 10^{-4}$. Quality factor $Q = 1/(2
\zeta) \approx 1060$. The assumption is that an audible
$40\,\si{\decibel}$ drop corresponds to a $1/100$ amplitude
ratio; the actual threshold of perception varies with frequency
and listener, and a different assumption shifts $Q$ by a factor
of two or three. The order of magnitude (a few thousand) is the
working number for steel-string guitar tones.
\end{exercise}

\begin{exercise}
Estimate the time scale on which the linear logistic model
$\dot{N} = r N (1 - N/K)$ transitions from approximately
exponential to approximately saturating, for a population whose
initial size is $N_{0} = K/100$ and growth rate $r =
0.1\,\si{\per\day}$.

\paragraph{Solution sketch.} The transition is the inflection
point, $N = K/2$. Starting from $N_{0} = K/100$ in the
exponential regime ($\dot{N} \approx r N$), the time to reach
$K/2$ is approximately the time for $50\times$ growth: $t \approx
\ln(K/2)/(r N_{0}) \cdot N_{0} = \ln(50)/r \approx 3.91/0.1
\approx 39\,\si{\day}$. More carefully, the logistic solution
$N(t) = K/(1 + ((K - N_{0})/N_{0}) e^{-r t})$ at $N = K/2$ gives
$(K - N_{0})/N_{0} \cdot e^{-r t} = 1$, hence $t = \ln(99) /
0.1 = 45.9\,\si{\day}$. Estimate and exact agree to a factor of
$1.2$. The transition is centred on $\sim 6$ weeks; before that
the trajectory looks exponential, after that it asymptotes to
$K$.
\end{exercise}

\begin{exercise}
Estimate the step size required for a forward-Euler integration
of the stiff system $\dot{y}_{1} = -1000 y_{1} + y_{2}$,
$\dot{y}_{2} = -y_{2}$ over a $1\,\si{\second}$ window. Compare
to the step size that an implicit Euler method would tolerate at
the same target accuracy.

\paragraph{Solution sketch.} The Jacobian is
$A = \begin{pmatrix} -1000 & 1 \\ 0 & -1 \end{pmatrix}$, upper
triangular with eigenvalues $\lambda_{1} = -1000$, $\lambda_{2} =
-1$. Forward Euler requires $h < 2/|\lambda_{1}| = 2 \times
10^{-3}\,\si{\second}$ for stability, set by the fast mode, even
after the fast mode has decayed below machine precision (which
happens at $t \approx 10/1000 = 10^{-2}\,\si{\second}$). Over the
one-second window, that mandates $N \geq 500$ steps for
stability alone; for a one-percent accuracy on the slow mode,
$h \sim 10^{-2}$ would suffice on accuracy grounds but stability
overrides. Take $h = 10^{-3}\,\si{\second}$, $N = 1000$ steps.
Implicit Euler is A-stable, so step size is set by accuracy on
the slow mode only: $h \sim 10^{-2}\,\si{\second}$ gives
percent-level error on $y_{2}(t) = y_{2}(0) e^{-t}$, $N = 100$
steps. Even allowing for the per-step linear solve (here cheap;
the system is $2 \times 2$), implicit Euler wins by an order of
magnitude in wall time. The stiffness ratio
$|\lambda_{1}|/|\lambda_{2}| = 1000$ is the working number: any
ratio above $\sim 100$ pushes the explicit method into the
penalty regime.
\end{exercise}

\begin{exercise}
Estimate the resonant amplification of a mechanical system whose
quality factor is $Q = 50$ when driven sinusoidally at exactly
its natural frequency. Estimate the steady-state amplitude when
the same system is driven at $1.5\omega_{0}$, well above
resonance. State the dimensionless ratio that controls the
high-frequency rolloff.

\paragraph{Solution sketch.} At resonance, the steady-state
amplitude is $A(\omega_{0}) \approx (F_{0}/k) Q = 50 \cdot
(F_{0}/k)$: a fifty-fold amplification of the static deflection.
Off resonance at $\omega = 1.5 \omega_{0}$, the denominator
$\sqrt{(k - m \omega^{2})^{2} + (c \omega)^{2}}$ is dominated by
the stiffness term: $k - m (1.5 \omega_{0})^{2} = k - 2.25 k =
-1.25 k$. The damping term $c \omega = c \cdot 1.5 \omega_{0} =
(2 \zeta k/\omega_{0})(1.5 \omega_{0}) = 3 \zeta k = 0.03 k$ (for
$\zeta = 1/100$), negligible compared to $1.25 k$. So
$A(1.5 \omega_{0}) \approx F_{0}/(1.25 k) = 0.8 (F_{0}/k)$. The
ratio $A(\omega_{0})/A(1.5 \omega_{0}) \approx 50/0.8 \approx 62$:
moving from on-resonance to $50\,\%$ above natural frequency
collapses the response by nearly two orders of magnitude. The
dimensionless ratio controlling the high-frequency rolloff is
$\omega/\omega_{0}$; the rolloff is $-40\,\si{\decibel}$ per
decade past resonance for a second-order system, set by the
$\omega^{-2}$ factor in $A(\omega)$ as $\omega \to \infty$.
\end{exercise}

\begin{exercise}
Estimate, before computing, the time at which a logistic
population $\dot{N} = 0.5 N (1 - N/1000)$ with $N(0) = 10$
reaches half of its carrying capacity. Compare to the formal
answer obtained by solving the closed-form expression.

\paragraph{Solution sketch.} Parameters: $r = 0.5\,\si{\per
\unit{time}}$, $K = 1000$, $N_{0} = 10$. The early dynamics are
approximately exponential, $N(t) \approx N_{0} e^{r t}$ until
$N \sim K$. Estimate the time for $N$ to grow from $10$ to $500$:
$500/10 = 50$, so $t \approx \ln 50 / r = 3.912 / 0.5 \approx
7.8\,\si{\unit{time}}$. The exact logistic solution is
\[
N(t) = \frac{K}{1 + ((K - N_{0})/N_{0})\, e^{-r t}}.
\]
Setting $N(t_{1/2}) = K/2 = 500$ gives $(K - N_{0})/N_{0} \cdot
e^{-r t_{1/2}} = 1$, hence $e^{r t_{1/2}} = (K - N_{0})/N_{0} =
990/10 = 99$. Then $t_{1/2} = \ln 99 / 0.5 = 4.595 / 0.5 \approx
9.19\,\si{\unit{time}}$. The estimate $7.8$ and the exact $9.19$
differ by $15\,\%$; the exponential approximation slightly
under-predicts because by $N = 500$ the logistic curve has
already begun to bend. The half-capacity time is the inflection
point of $N(t)$, the moment at which $\dot{N}$ peaks at $r K/4 =
125$.
\end{exercise}

\subsection*{Simulation}\pathtag{standard}

\begin{exercise}
Implement forward Euler and RK4 in a programming environment of
your choice. Apply each to the test problem $\dot{y} = -y$,
$y(0) = 1$, on the interval $[0, 5]$. Plot the absolute error at
$t = 5$ as a function of step size for $h \in \{0.5, 0.1, 0.01,
0.001\}$. Confirm the predicted slopes ($1$ for Euler, $4$ for
RK4) on log-log axes.

\paragraph{Solution sketch.} The exact solution at $t = 5$ is
$y(5) = e^{-5} \approx 6.738 \times 10^{-3}$. Forward Euler
applied to $\dot{y} = -y$ gives $y_{n+1} = (1 - h) y_{n}$, so
after $N = 5/h$ steps the numerical value is $y_{N} = (1 -
h)^{5/h}$. Tabulate:
\[
\begin{array}{c|c|c|c}
h & y_{\text{Euler}}(5) & |\text{error}| & |\text{err}|/h \\\hline
0.5   & (0.5)^{10}   = 9.77 \times 10^{-4}   & 5.76 \times 10^{-3} & 1.15 \times 10^{-2} \\
0.1   & (0.9)^{50}   = 5.15 \times 10^{-3}   & 1.59 \times 10^{-3} & 1.59 \times 10^{-2} \\
0.01  & (0.99)^{500} = 6.57 \times 10^{-3}   & 1.69 \times 10^{-4} & 1.69 \times 10^{-2} \\
0.001 & (0.999)^{5000} = 6.72 \times 10^{-3} & 1.68 \times 10^{-5} & 1.68 \times 10^{-2}
\end{array}
\]
The error-to-step-size ratio stabilises near $1.7 \times 10^{-2}$
as $h \to 0$: the leading-order constant in $|\text{error}|
\approx C h$ is approximately $C = (T/2) |y''(\xi)| \approx (5/2)
\cdot e^{-\xi}$ for $\xi \in [0, 5]$, dominated by the early
interval where $y$ is large; numerically $C \approx 0.017$. On
log-log axes, the four points fall on a line of slope $1$.

For RK4, the amplification factor per step on $\dot{y} = -y$ is
$R(h) = 1 - h + h^{2}/2 - h^{3}/6 + h^{4}/24$, the truncated
exponential. After $N$ steps, $y_{N} = R(h)^{N}$. Tabulate:
\[
\begin{array}{c|c|c|c}
h & y_{\text{RK4}}(5) & |\text{error}| & |\text{err}|/h^{4} \\\hline
0.5   & 6.79 \times 10^{-3}      & 5.1 \times 10^{-5}    & 8.1 \times 10^{-4} \\
0.1   & 6.7382 \times 10^{-3}    & 1.39 \times 10^{-7}   & 1.39 \times 10^{-3} \\
0.01  & 6.7379469 \times 10^{-3} & 1.41 \times 10^{-11}  & 1.41 \times 10^{-3} \\
0.001 & \text{at machine eps}    & \approx 10^{-15}      & \text{floor}
\end{array}
\]
The ratio $|\text{err}|/h^{4}$ stabilises near $1.4 \times
10^{-3}$ for the middle two values; on log-log axes the slope is
exactly $4$ between $h = 0.5$ and $h = 0.01$. At $h = 0.001$ the
truncation error is below double-precision representation and
the error floor is set by accumulated roundoff, $\sim N
\epsilon \approx 5000 \cdot 2.2 \times 10^{-16} \approx 10^{-12}$:
the slope flattens at the precision floor, the classical signature
that RK4 at small $h$ on a benign problem is limited by
arithmetic rather than by truncation. The reader who confirms the
two slopes by numerical experiment has confirmed both the order
of accuracy theorem and the roundoff floor that bounds it from
below; see \cite{text:v2ch07-hairer-norsett-wanner} for the
canonical treatment.
\end{exercise}

\begin{exercise}
Simulate the underdamped oscillator $\ddot{x} + 0.1 \dot{x} + x =
0$ with $x(0) = 1$, $\dot{x}(0) = 0$. Plot $x(t)$ and the
amplitude envelope $e^{-\zeta\omega_{0} t}$ on the same axes.
Extract the damped frequency from the simulation by zero-crossing
detection and compare to the analytical $\omega_{d}$.

\paragraph{Solution sketch.} Identify $m = 1$, $c = 0.1$, $k = 1$:
natural frequency $\omega_{0} = 1\,\si{\radian\per\second}$,
damping ratio $\zeta = 0.05$, damped frequency $\omega_{d} =
\omega_{0}\sqrt{1 - \zeta^{2}} = \sqrt{0.9975} \approx 0.99875\,
\si{\radian\per\second}$. Decay rate $\zeta \omega_{0} = 0.05$.
Closed-form solution: $x(t) = e^{-0.05 t}(\cos\omega_{d} t +
(0.05/\omega_{d}) \sin\omega_{d} t) \approx e^{-0.05 t} \cos
\omega_{d} t$ (the sine term has small coefficient). Convert to
the first-order system $\mathbf{y} = (x, \dot{x})^{T}$,
$\dot{\mathbf{y}} = (\dot{x}, -x - 0.1 \dot{x})^{T}$ and integrate
with RK4 at $h = 0.01$ over $[0, 100]$. The amplitude envelope
$e^{-0.05 t}$ falls to $1/e$ at $t = 20$, to $1\,\%$ at $t = 92$.
Period $T_{d} = 2\pi/\omega_{d} \approx 6.291\,\si{\second}$;
about $16$ periods fit in the run. Zero-crossing extraction:
identify all $n$ where $\operatorname{sgn}(x_{n}) \neq
\operatorname{sgn}(x_{n-1})$; interpolate linearly between
$(t_{n-1}, x_{n-1})$ and $(t_{n}, x_{n})$ for sub-step
resolution. Spacing between successive upward (negative-to-
positive) crossings gives $T_{d}$. Numerical: $T_{d, \text{num}}
\approx 6.291\,\si{\second}$, matching analytic to $5 \times
10^{-5}$. The envelope sampled at every upward crossing should
fit $\ln |x_{\text{peak}}| = -\zeta \omega_{0} t + $ const with
slope $-0.05$; the slope extracted by least squares gives
$\zeta \omega_{0}$ directly, the standard logarithmic-decrement
method for damping identification.
\end{exercise}

\begin{exercise}
Simulate the forced oscillator $\ddot{x} + 0.1 \dot{x} + x =
\cos(\omega t)$ for $\omega \in \{0.5, 0.9, 1.0, 1.1, 1.5\}$.
Allow each simulation to run long enough for the transient to
decay (at least $100$ time units). Extract the steady-state
amplitude and phase. Plot the amplitude and phase as functions of
$\omega$ and compare to the analytical formulas.

\paragraph{Solution sketch.} With $m = k = 1$, $c = 0.1$, $F_{0}
= 1$: $\omega_{0} = 1$, $\zeta = 0.05$, $Q = 10$. Analytical
amplitude $A(\omega) = 1/\sqrt{(1 - \omega^{2})^{2} + (0.1
\omega)^{2}}$; analytical phase lag $\varphi(\omega) = \arctan
(0.1 \omega / (1 - \omega^{2}))$, with $\varphi \in [0, \pi]$
(branch chosen so that $\varphi$ is continuous and equals
$\pi/2$ at $\omega = 1$). Tabulate:
\[
\begin{array}{c|c|c}
\omega & A(\omega) & \varphi(\omega) \\\hline
0.5 & 1/\sqrt{0.5625 + 0.0025} = 1.331 & 0.067\,\si{\radian} \\
0.9 & 1/\sqrt{0.0361 + 0.0081} = 4.756 & 0.444\,\si{\radian} \\
1.0 & 1/0.1 = 10.00 & \pi/2 = 1.571\,\si{\radian} \\
1.1 & 1/\sqrt{0.0441 + 0.0121} = 4.217 & \pi - 0.488 = 2.654\,\si{\radian} \\
1.5 & 1/\sqrt{1.5625 + 0.0225} = 0.794 & \pi - 0.119 = 3.023\,\si{\radian}
\end{array}
\]
The peak amplitude $A(1.0) = Q = 10$ is the resonance maximum.
Simulation: integrate from $x(0) = \dot{x}(0) = 0$ over $[0,
200]$ with RK4, $h = 0.01$. The transient envelope decays as
$e^{-0.05 t}$, falling to $1\,\%$ near $t = 92$; truncate the
first $100$ time units and analyse the tail. Steady-state
amplitude: fit $x(t) = A \cos(\omega t - \varphi)$ to the last
$100$ time units. Numerical amplitude and phase agree with the
analytical values to within $0.5\,\%$ at all five frequencies,
provided the transient has decayed. Common mistake: truncating
too early. If the simulation starts the analysis at $t = 50$,
the residual transient amplitude $e^{-2.5} \approx 0.082$
biases the amplitude estimate by $\sim 1\,\%$ at off-resonance
frequencies and more near resonance where the transient and
steady-state interfere. The lesson: settling time
$\sim Q/\omega_{0} = 10$ time constants of the slow damping,
$\sim 100\,\si{\second}$ for $Q = 10$; respect this.
\end{exercise}

\begin{exercise}
Simulate the predator-prey system $\dot{x} = x - 0.5 x y$,
$\dot{y} = -0.75 y + 0.25 x y$ with $x(0) = 5$, $y(0) = 2$, over
$1{,}000$ time units. Plot the trajectory in the $(x, y)$ plane.
Compute the conserved quantity $H(x, y)$ along the trajectory and
plot it as a function of time; observe the drift. Repeat with a
symplectic integrator and compare.

\paragraph{Solution sketch.} Parameters $a = 1$, $b = 0.5$, $c =
0.75$, $d = 0.25$. Coexistence equilibrium $(c/d, a/b) = (3, 2)$
(linearisation gives a centre with predicted period $T = 2\pi/
\sqrt{a c} \approx 7.26\,\si{\unit{time}}$). The conserved
quantity is $H(x, y) = 0.25 x - 0.75 \ln x + 0.5 y - \ln y$.
Evaluate at the initial point: $H(5, 2) = 1.25 - 0.75 \ln 5 + 1
- \ln 2 = 1.25 - 1.207 + 1 - 0.693 = 0.350$. RK4 integration
with $h = 0.01$ over $[0, 1000]$ produces a trajectory that
traces a closed orbit for the first few hundred time units;
$H(t)$ along the trajectory drifts upward (or downward,
depending on the local phase-space curvature) by a few percent
over the full run. Cause: RK4's modified equation has a
non-zero contribution that is not Hamiltonian; over many
periods, $H$ accumulates a secular drift. Repeating with a
second-order symplectic integrator (e.g.\ St\"ormer-Verlet
applied after a Poincar\'e transformation to canonical
coordinates, or simply the symplectic Euler $x_{n+1} = x_{n} +
h f_{1}(x_{n}, y_{n})$, $y_{n+1} = y_{n} + h f_{2}(x_{n+1},
y_{n})$) shows $H$ oscillating with bounded amplitude
proportional to $h^{2}$, with no secular drift. The orbit
remains closed to within machine precision for the full run.
The exercise is the working demonstration of the symplectic
discipline: the right integrator for the right problem matters
more than the order of accuracy. See
\cite{text:v2ch07-hairer-lubich-wanner-geometric} for the full
theory.
\end{exercise}

\begin{exercise}
Simulate the stiff system $\dot{y}_{1} = -1000 y_{1} + 999 y_{2}$,
$\dot{y}_{2} = y_{1} - 2 y_{2}$ with $y_{1}(0) = 1$, $y_{2}(0) =
0$, on the interval $[0, 10]$ using forward Euler with $h =
0.001$, $0.0015$, $0.0025$. Identify the step size at which the
explicit method becomes unstable. Solve the same problem with
implicit Euler at $h = 0.1$ and confirm the slow mode is
captured correctly.

\paragraph{Solution sketch.} Jacobian
$A = \begin{pmatrix} -1000 & 999 \\ 1 & -2 \end{pmatrix}$.
Trace $T = -1002$, determinant $D = 2000 - 999 = 1001$.
Eigenvalues: $\lambda = (T \pm \sqrt{T^{2} - 4 D})/2 = (-1002
\pm \sqrt{1{,}004{,}004 - 4004})/2 = (-1002 \pm \sqrt{1{,}000{,}000})/2
= (-1002 \pm 1000)/2$, giving $\lambda_{1} = -1$ (slow) and
$\lambda_{2} = -1001$ (fast). Stiffness ratio $\approx 1000$.
Forward-Euler stability requires $h < 2/|\lambda_{2}| = 2/1001
\approx 0.001998\,\si{\second}$. Predictions: $h = 0.001$ is
stable (amplification factor $|1 - 1.001| = 0.001 < 1$ for the
fast mode); $h = 0.0015$ is stable but marginally so ($|1 -
1.5015| = 0.5015 < 1$); $h = 0.0025$ exceeds the bound,
amplification factor $|1 - 2.5025| = 1.5025 > 1$, the fast mode
amplifies geometrically and the solution blows up within tens of
steps. The numerical experiment confirms exactly this: $h =
0.001$ and $h = 0.0015$ produce smooth decaying trajectories
matching the analytical solution to four decimal places at
$t = 10$; $h = 0.0025$ produces a runaway oscillation, with
$y_{1}$ alternating sign at the Nyquist frequency $1/(2 h)$ and
growing by a factor of $1.5$ per step. Implicit Euler with
$h = 0.1$ requires solving the linear system $(I - h A) y_{n+1}
= y_{n}$ each step, but that solve is one $2 \times 2$ inverse
($\sim 10$ flops) and produces a stable, accurate slow-mode
trajectory in $100$ steps versus $10{,}000$ for the stable
explicit run. The slow-mode time constant is $1/|\lambda_{1}|
= 1\,\si{\second}$; the solution at $t = 10$ has decayed to
$\sim 10^{-5}$ of its initial value (slow mode decayed by
$e^{-10} \approx 4.5 \times 10^{-5}$). The reader observes the
qualitative difference: stiffness reflects a mismatch between
method and problem, not difficulty of the underlying physics.
\end{exercise}

\begin{exercise}
Simulate the simple pendulum $\ddot\theta + \sin\theta = 0$ with
$\theta(0) = 1.0\,\si{\radian}$ ($\approx 57^{\circ}$),
$\dot\theta(0) = 0$. Use RK4 with $h = 0.01$ over $50$ time
units. Extract the period and compare to the small-angle
prediction $T_{0} = 2\pi$ and the first-correction prediction
$T_{1} = T_{0}(1 + \theta_{0}^{2}/16)$. Compute the energy drift
over the run.

\paragraph{Solution sketch.} Set $\omega_{0}^{2} = 1$, so the
linear period is $T_{0} = 2 \pi \approx 6.2832$ time units.
First correction: $T_{1} = T_{0}(1 + 1.0^{2}/16) = T_{0} \cdot
1.0625 = 6.6759$. The exact period from the elliptic integral
of the first kind is $T_{\text{ex}} = 4 K(\sin(\theta_{0}/2)) =
4 K(\sin 0.5) = 4 K(0.4794) \approx 6.6745$ (using tabulated
values or numerical quadrature of $K$). First correction matches
to $\sim 0.02\,\%$; the small-angle prediction is short by
$\sim 6.3\,\%$. Convert to first-order: $\mathbf{y} = (\theta,
\dot\theta)^{T}$, $\dot{\mathbf{y}} = (\dot\theta, -\sin\theta)^{T}$.
RK4 with $h = 0.01$: about $5000$ steps. Period extraction by
zero-crossing of $\dot\theta$ (positive-to-negative, at the top
of the swing) averaged over the $\sim 7$ to $8$ oscillations in
the run gives $T_{\text{num}} \approx 6.6745 \pm 10^{-4}$,
matching the exact value. Energy: $E(t) = (1/2) \dot\theta^{2}
+ (1 - \cos\theta)$. Initial value $E(0) = 1 - \cos 1.0 \approx
0.4597$. Over the $50$-time-unit run, RK4 produces $|\Delta E|/E
\sim 10^{-6}$ to $10^{-5}$, a slow drift from the
non-symplectic modified equation; over longer runs the drift
accumulates linearly. A symplectic leapfrog at the same $h$
would keep $|\Delta E|/E$ oscillating with bounded amplitude $\sim
h^{2} = 10^{-4}$ but no secular trend. The exercise is the
working bridge to the chapter project, which sweeps amplitude
from $5^{\circ}$ to $90^{\circ}$ and quantifies the
small-angle-vs-first-correction-vs-numerical comparison across
the range.
\end{exercise}

\subsection*{Design}\pathtag{standard}

\begin{exercise}
Design an RC low-pass filter with $-3\,\si{\decibel}$ corner
frequency at $1\,\si{\kilo\hertz}$ for an audio signal-conditioning
chain. State the values of $R$ and $C$, the transient response
time, and the filter's behaviour at $f = 100\,\si{\hertz}$ and
$f = 10\,\si{\kilo\hertz}$. Identify a regime where the
first-order rolloff is not sharp enough and a second-order
filter would be required.

\paragraph{Discussion.} The $-3\,\si{\decibel}$ corner of the RC
low-pass is $f_{c} = 1/(2 \pi R C) = 10^{3}\,\si{\hertz}$, giving
$R C = 1/(2 \pi \times 10^{3}) \approx 1.59 \times 10^{-4}\,
\si{\second}$. A well-formed answer picks one degree of freedom
from the working component palette: e.g., $C = 10\,\si{\nano
\farad}$, hence $R = 15.9\,\si{\kilo\ohm}$ (or any equivalent
pair from the $E12$ or $E24$ series). Transient response time
constant $\tau = R C = 1.59 \times 10^{-4}\,\si{\second}$; a step
input reaches $63\,\%$ in one $\tau$, $95\,\%$ in three $\tau$,
$99.9\,\%$ in seven $\tau$ ($\sim 1.1\,\si{\milli\second}$). At
$f = 100\,\si{\hertz}$ ($f/f_{c} = 0.1$), the magnitude is
$1/\sqrt{1 + 0.01} \approx 0.995$, attenuation $\sim 0.04\,
\si{\decibel}$, phase lag $\arctan 0.1 \approx 5.7^{\circ}$. At
$f = 10\,\si{\kilo\hertz}$ ($f/f_{c} = 10$), magnitude $1/
\sqrt{1 + 100} \approx 0.0995$, attenuation $-20\,\si{\decibel}$,
phase lag $\arctan 10 \approx 84.3^{\circ}$. Failure mode: the
first-order rolloff at $-20\,\si{\decibel}$/decade is too gentle
for anti-aliasing in front of a sampled-data system whose
Nyquist frequency is close to $f_{c}$; the readout of frequencies
above the sample rate aliases in. The standard correction is a
higher-order filter (Sallen-Key second-order or
biquad/multi-feedback structures) with $-40\,\si{\decibel}$ or
steeper rolloff. The editor would accept any pair $(R, C)$ that
sets $R C$ to the right value, any reasonable transient-time
statement, and any concrete identification of an over-stretched
single-pole filter (anti-aliasing, mains-hum rejection on a
weak signal, etc.).
\end{exercise}

\begin{exercise}
Design a critically damped door closer for a $5\,\si{\kilogram}$
door. Choose the spring constant $k$ to give the desired closing
time of approximately $4\,\si{\second}$ for a 90-degree opening
and the damping coefficient $c$ to satisfy the
critical-damping condition $c = 2\sqrt{m k}$. State the
assumptions about the moment of inertia and the dependency on
door geometry; identify which assumption dominates the design
uncertainty.

\paragraph{Discussion.} The well-formed answer recognises that
the door rotates about a hinge axis, so the working dynamics are
$I \ddot\theta + c \dot\theta + k \theta = 0$ with moment of
inertia $I$, not mass $m$. For a rectangular door of mass
$5\,\si{\kilogram}$, width $w = 0.8\,\si{\meter}$, with the
hinge axis at one edge, $I = (1/3) m w^{2} = (1/3) \cdot 5 \cdot
0.64 \approx 1.07\,\si{\kilogram\meter\squared}$. Critical
damping sets $\zeta = 1$: the response to a $90^{\circ}$ initial
displacement is $\theta(t) = \theta_{0}(1 + \omega_{0} t) e^{-
\omega_{0} t}$. The $99\,\%$-closed time is approximately $7/
\omega_{0}$ (the $(1 + \omega_{0} t)$ envelope decays at $1\,\%$
of $\theta_{0}$ near $\omega_{0} t \approx 7$). A four-second
target gives $\omega_{0} \approx 7/4 = 1.75\,\si{\radian\per
\second}$, hence $k = I \omega_{0}^{2} \approx 1.07 \cdot 3.06
\approx 3.3\,\si{\newton\meter\per\radian}$ (a torsional spring
constant; the working unit). Damping coefficient $c = 2 \sqrt{I
k} \approx 2 \sqrt{3.5} \approx 3.7\,\si{\newton\meter\second
\per\radian}$. Failure modes the editor watches for: confusing
$m \ddot{x} + c \dot{x} + k x = 0$ (translational) with the
rotational form; ignoring the moment of inertia; mis-defining
the $99\,\%$ closing time (a critically damped response has no
overshoot but also no fast decay). The dominant design
uncertainty is the door geometry, specifically the width and the
mass distribution; the moment of inertia varies as $w^{2}$ at
fixed mass, so a $20\,\%$ width error propagates to $44\,\%$ in
$I$ and to a $20\,\%$ shift in $\omega_{0}$, hence the closing
time. The editor accepts answers that name geometry as the
dominant source and that arrive at a working $k$, $c$ pair
through the rotational form.
\end{exercise}

\begin{exercise}
Design a numerical integration strategy for a long-running
orbital simulation: a satellite around Earth tracked over one
year. State the integrator family (symplectic versus
non-symplectic), the step size, the tolerance, the conserved
quantity used for runtime monitoring, and the threshold at which
the simulation triggers a manual review. Identify the dominant
expected source of long-term integration error.

\paragraph{Discussion.} The well-formed answer picks a symplectic
or geometric integrator and justifies the choice on long-term
energy and angular-momentum behaviour. A working choice: a
fourth-order symplectic integrator such as the Forest-Ruth or
Yoshida 4th-order method (a product of symplectic Euler
sub-steps), or a Wisdom-Holman mixed-variable scheme adapted to
near-Keplerian orbits. Step size for a low-Earth orbit of period
$\sim 90\,\si{\minute}$: $h = 30\,\si{\second}$ resolves one
orbit in $\sim 180$ steps, giving phase-angle accuracy of
$\sim 10^{-4}$ per orbit at fourth order. Over a one-year run
($\sim 5800$ orbits, $\sim 10^{6}$ steps), the symplectic
modified equation keeps $|\Delta E|/E$ oscillating with bounded
amplitude $\sim h^{4}$; angular momentum is preserved to
machine precision for an axisymmetric potential. Conserved
quantity for runtime monitoring: the Kepler invariant $E =
(1/2)|\dot{\mathbf{r}}|^{2} - GM/|\mathbf{r}|$ at every step
(cheap; one dot product and one inverse-norm), with a sliding
window of $1000$ steps tracked for trend and amplitude. The
review trigger: a secular trend in $|\Delta E|/E$ exceeding
$10^{-8}$ per orbit over a $100$-orbit window, which would
signal either a step-size violation in a high-eccentricity
passage (close approach to Earth) or a numerical-precision
floor that has been crossed. Dominant long-term error source:
roundoff accumulation, which grows as $\sqrt{N} \epsilon$ at
double precision for a symplectic method, $\sim \sqrt{10^{6}}
\cdot 2.2 \times 10^{-16} \approx 2 \times 10^{-13}$ per phase
variable, well below the orbit-determination tolerance.
Secondary: the non-Keplerian perturbations (Earth oblateness
$J_{2}$, atmospheric drag at low altitude, lunar and solar
gravity) the model includes or omits; the integrator error is
a small fraction of the modelling error for any realistic
satellite. Failure mode the editor watches for: choosing an
explicit non-symplectic method with adaptive step size, which
introduces uncontrolled energy drift at every accepted-step
event. Acceptable answer: any symplectic-class integrator with a
defended step size, an explicit conservation monitor, and a
named trigger threshold.
\end{exercise}

\begin{exercise}
Design a dosing schedule for a one-compartment drug whose
half-life is $4\,\si{\hour}$ that maintains the trough
concentration above a minimum effective concentration $C_{\min}$
and the peak concentration below a toxicity threshold $C_{\max}$
in a steady-state regimen. State the dose, the dosing interval,
and the time to reach steady state. Identify the design
parameter that is most sensitive to inter-patient variability
in the elimination rate constant.

\paragraph{Discussion.} Elimination constant $k = \ln 2 / 4 =
0.173\,\si{\per\hour}$. At steady state with dose $D$ and
interval $\tau$, the trough and peak concentrations are
$C_{\text{trough}} = D/V \cdot e^{-k \tau}/(1 - e^{-k \tau})$
and $C_{\text{peak}} = D/V \cdot 1/(1 - e^{-k \tau})$ (the
geometric-series sum of decaying boluses). The well-formed
answer derives the ratio $C_{\text{peak}}/C_{\text{trough}} =
e^{k \tau}$, sets it equal to the therapeutic-to-trough ratio
$C_{\max}/C_{\min}$, and solves for $\tau$: $\tau = \ln
(C_{\max}/C_{\min}) / k$. Worked example with a $5\times$
therapeutic window: $\tau = \ln 5 / 0.173 \approx 9.3\,\si{\hour}$,
rounded operationally to $8\,\si{\hour}$ (three-times-daily) or
$12\,\si{\hour}$. With $\tau = 8\,\si{\hour}$, $e^{-k \tau} =
e^{-1.39} = 0.249$, so $C_{\text{trough}} = (D/V) \cdot 0.249 /
0.751 = (D/V) \cdot 0.332$. Choose $D$ so $C_{\text{trough}} =
C_{\min}$: $D/V = 3.0 C_{\min}$, giving $C_{\text{peak}} =
3.0 C_{\min} / 0.751 \approx 4.0 C_{\min}$, safely below
$5 C_{\min}$. Time to steady state: about $4$ half-lives
($\sim 16\,\si{\hour}$, two dosing intervals), to within $6\,\%$
of the steady-state values. Sensitivity: trough concentration
is the most sensitive parameter. $\partial \ln C_{\text{trough}}
/\partial \ln k = -k \tau / (1 - e^{-k \tau}) \approx -1.85$ for
$k \tau \approx 1.39$. A $30\,\%$ slow elimination (clearance
reduced by $30\,\%$) drives the trough up by $\sim 55\,\%$,
potentially into toxic territory; a $30\,\%$ fast elimination
drives the trough down by $\sim 55\,\%$, potentially
sub-therapeutic. The clinical correction is a loading dose plus
therapeutic-drug monitoring for narrow-therapeutic-window agents
(aminoglycosides, vancomycin, lithium). The editor accepts any
answer that derives the peak/trough ratio formula, identifies
$\tau$ as the dose-spacing knob, and names the trough as the
inter-patient-variability danger.
\end{exercise}

\subsection*{Diagnosis and reverse engineering}\pathtag{core}

\begin{exercise}
A laboratory ODE simulation of a cooling object's temperature
shows the temperature continuously oscillating with growing
amplitude despite the physics predicting monotone decay. State
three plausible mechanisms (forward Euler stability violation,
sign error in the right-hand side, time-step-dependent driving
term) and propose the diagnostic test that distinguishes them.

\paragraph{Solution sketch.} Mechanism 1 (forward Euler
stability): the time constant of the cooling is $\tau = 1/k$,
and the explicit-Euler stability bound is $h < 2 \tau$. A step
size near or above $2 \tau$ produces the diagnostic signature of
period-$2 h$ alternation with growth factor $|1 - h k| > 1$.
Diagnostic test: halve the step size and rerun. If the
oscillation disappears (with the new amplification factor closer
to $1$), the cause is stability. Mechanism 2 (sign error in the
right-hand side): the user has coded $\dot{T} = +k(T - T_{\text
{env}})$ instead of $\dot{T} = -k(T - T_{\text{env}})$, giving
pure exponential growth, not oscillation. Diagnostic: inspect
the code's sign and run a short integration analytically; pure
exponential divergence with no sign change rules out
oscillatory instability. Mechanism 3 (time-step-dependent
driving): an explicit driving term computed at $t_{n}$ but
intended to be at $t_{n} + h$ introduces a phase lag that, in
combination with a feedback loop (e.g., an integrated controller
state), produces a discrete-time oscillation whose frequency
locks to $1/h$. Diagnostic: rerun at the same $h$ with the
driving evaluated at $t_{n + 1/2}$ (midpoint rule); if the
oscillation halves in amplitude, the driving was the cause. The
discipline: each diagnostic isolates a single mechanism, and the
test is cheap. The wrong response is to "tune" the step size by
trial and error until the symptom disappears; the right response
is to halve, observe, and infer.
\end{exercise}

\begin{exercise}
A simulation of a chemical reactor's concentration uses an
explicit RK4 integrator with $h = 10^{-3}\,\si{\second}$. The
simulation runs for several minutes of model time, then the
output blows up. The reactor's fastest mode has time constant
$\sim 10^{-4}\,\si{\second}$. State the most likely cause and
the corrective action.

\paragraph{Solution sketch.} The fast time constant $\tau_{f}
\sim 10^{-4}\,\si{\second}$ corresponds to $|\lambda_{f}| \sim
10^{4}\,\si{\per\second}$. RK4 stability for $\dot{y} = \lambda
y$ requires $h |\lambda| < 2.78$ (the rightmost point of the
RK4 stability region on the negative real axis); at $h =
10^{-3}$, $h |\lambda_{f}| = 10$, well outside the stability
region. The fast mode amplifies by a factor of $|R(-10)|$
per step, where $R$ is the RK4 stability function; $|R(-10)|
\approx 9.8$, so a single Nyquist-frequency oscillation grows
$\sim 9.8\times$ per step. After enough seconds of model time
for the accumulated fast-mode noise to dominate, the output
explodes. The fact that several minutes elapse before the
blow-up is the diagnostic clue: the slow modes are tracked
correctly, the fast mode is suppressed at machine-precision
amplitude initially, and the runaway sets in only once the
amplified noise has grown to a level that perturbs the slow
modes via nonlinear coupling. Corrective action: switch to an
implicit method (BDF or implicit Runge-Kutta or a Rosenbrock
scheme) and select step size by accuracy on the slow modes,
e.g., $h = 10^{-2}\,\si{\second}$. Alternative: keep RK4 but
reduce step size below $10^{-4} \cdot 2.78 \approx 2.8 \times
10^{-4}\,\si{\second}$, an order of magnitude smaller; this
works but costs $\sim 4\times$ more wall time than the implicit
solution on a stiff problem.
\end{exercise}

\begin{exercise}
A pharmacokinetic simulation of a multi-dose regimen shows
steady-state drug concentrations $50\,\%$ above the predictions
of a one-compartment linear model. State three plausible
mechanisms (saturation of elimination kinetics, accumulation due
to inadequate clearance interval, second compartment with slower
exchange) and propose the experimental measurement that would
distinguish them.

\paragraph{Solution sketch.} Mechanism 1 (Michaelis-Menten
saturation): at high concentrations the clearance plateaus at
$V_{\max}$, so the effective half-life lengthens with dose. The
elevation grows with dose; the steady-state error vanishes at
low doses. Diagnostic: a dose-ranging study, with the offset
quantified at several dose levels. If the $50\,\%$ elevation
grows with dose, saturation is implicated. Mechanism 2
(inadequate clearance interval): if the dosing interval is
shorter than the steady-state assumption requires (e.g., the
study used $\tau < 4 t_{1/2}$ to reach steady state but the
actual half-life is longer than assumed), the drug accumulates
beyond the model's prediction. Diagnostic: measure the
elimination rate constant directly from a single-dose washout
curve and compare to the value used in the simulation. If the
true $k$ is lower than assumed, the working half-life is
longer, and the accumulation prediction was wrong. Mechanism 3
(second compartment): a peripheral tissue compartment with
slow exchange acts as a reservoir, returning drug to plasma
over a long time constant. Steady-state concentration in the
central compartment is elevated because the peripheral
compartment buffers the elimination. Diagnostic: high-density
sampling of the post-dose concentration-time curve, with
biexponential fitting; a second compartment shows a long
elimination tail invisible in a single-exponential fit. The
working pharmacokinetic discipline: dose-ranging plus
multi-time-point sampling distinguishes the three on a single
study, and the choice of subsequent model (Michaelis-Menten,
revised one-compartment, or two-compartment) follows from the
diagnostic outcome.
\end{exercise}

\begin{exercise}
A predator-prey simulation produces trajectories that slowly
spiral into the coexistence equilibrium over $1{,}000$ time
units, contrary to the textbook prediction of closed orbits.
State the most likely cause (numerical dissipation in the
integrator) and propose the measurement that confirms it
(comparison against a symplectic integrator). State a second
candidate (model includes a small dissipation term that the
analyst missed) and propose the verification.

\paragraph{Solution sketch.} The textbook Lotka-Volterra system
has the centre at $(c/d, a/b)$ as a structurally fragile
equilibrium: any small perturbation (numerical or physical)
breaks the closed-orbit structure. Cause 1 (numerical
dissipation): a non-symplectic integrator (RK4, adaptive
Cash-Karp, classical Adams-Bashforth) has a modified equation
that is not Hamiltonian, and over many oscillations the
solution spirals slowly inward (or outward, depending on the
sign of the leading non-Hamiltonian term). Measurement to
confirm: rerun with a symplectic integrator (St\"ormer-Verlet
or symplectic Euler in canonical coordinates) at the same step
size; if the spiral disappears, the original integrator was the
cause. Quantitative check: $H(x, y) = d x - c \ln x + b y - a
\ln y$ along the trajectory should be constant; with the
non-symplectic method it drifts secularly, with the symplectic
method it oscillates with bounded amplitude $\sim h^{p}$ at
order $p$. Cause 2 (model error): the simulated right-hand side
includes a small dissipation term that the analyst missed:
$\dot{x} = a x - b x y - \epsilon x^{2}$, a logistic
correction with $\epsilon \ll 1$, which the textbook centre
does not contain. Verification: read the code carefully, set
$\epsilon = 0$ explicitly, and rerun; if the spiral persists,
the integrator is the cause; if it disappears, the model term
was the cause. The discipline transferring to other modelling
problems: any qualitative discrepancy between simulation and
textbook prediction has two candidate sources (integrator,
model), and the analyst's job is to isolate which one.
\end{exercise}

\subsection*{Failure analysis}\pathtag{core}

\begin{exercise}
The Patriot case at Dhahran (1991) is canonical for accumulated
arithmetic error in long-running numerical computation. Argue, in
200 to 300 words, what additional verification or operational
discipline would have detected the mechanism before the Dhahran
incident, citing the GAO report \cite{acc:gao-patriot-1992}.
Identify the analogue in long-running ODE integration: the
runtime estimate of accumulated roundoff that modern simulation
codes maintain.

\paragraph{Discussion.} The well-formed answer names three
verification steps that would have surfaced the
$0.34\,\si{\second}$ accumulated error before deployment to
Dhahran. First, a worst-case end-to-end timing test running the
system continuously for a $100$-hour interval, comparing
predicted range-gate positions against a known-truth ground
target track at one-hour intervals, would have exposed the
linear drift in clock conversion. The GAO report explicitly
notes that the error had been observed in field tests; the
verification gap was that "observed" did not propagate to
"systematically characterised and bounded"
\cite{acc:gao-patriot-1992}. Second, a static analysis of the
$24$-bit fixed-point representation of $1/10$ would have
identified the per-conversion error of $\approx 9.5 \times 10^{
-8}\,\si{\second}$ and the linear-in-$N$ accumulation pattern;
a back-of-envelope multiplication ($3.6 \times 10^{6} \times
9.5 \times 10^{-8}$) shows the $0.34\,\si{\second}$ figure.
Third, an operational discipline of mandatory restart after
$N$-hour runtime, tied to the documented characterisation,
would have made the error operationally invisible. The
analogue in ODE integration is the runtime estimate of
accumulated roundoff: modern simulation codes carry, alongside
the integrated state, a scalar estimate of cumulative
floating-point error growth (often a deterministic upper bound
$N \epsilon \|y\|$ and a probabilistic estimate $\sqrt{N}
\epsilon \|y\|$); the simulation halts or switches to extended
precision when the estimate exceeds a configured tolerance.
Long-running orbital mechanics codes, climate models, and
plasma simulations carry this discipline routinely; the
discipline is the numerical analogue of the operational
restart that Patriot's design had assumed and the deployment
had discarded.
\end{exercise}

\begin{exercise}
A spacecraft attitude-control simulation uses a non-symplectic
integrator and accumulates a small angular-momentum drift over
a multi-year mission. Argue, in 200 to 300 words, the failure
mode by which the drift could compromise mission performance,
the diagnostic that would detect it during testing, and the
discipline (symplectic integrator, runtime conservation
monitoring, periodic recalibration against an independent
reference) that bounds the failure.

\paragraph{Discussion.} The failure mode propagates from a
simulation artefact to a flight-mission consequence in three
stages. First, the non-symplectic integrator's modified equation
includes a small non-conservative component proportional to a
power of the step size; over $10^{8}$ steps in a multi-year
mission, even a $10^{-10}$ per-step drift in angular momentum
$|L|$ accumulates to a $10^{-2}$ fractional error. Second, the
flight controller, trained or tuned on the simulation, develops
an offset between commanded and actual reaction-wheel speeds
because the simulated dynamics imply a different momentum
balance than the actual ones; the offset manifests as a slow
bias in pointing accuracy. Third, the bias either drifts into
the science-instrument boresight tolerance (lost data) or
forces a momentum dump that exhausts attitude-control propellant
ahead of schedule (early end of mission). Diagnostic during
testing: a long-duration simulation run (a year of model time)
with the conserved quantity $|L|$ tracked at every step; a
plot of $|L|(t)$ versus $t$ shows either a bounded oscillation
(acceptable, with the integrator's modified equation still
near-conservative) or a secular trend (failure). A linear-fit
slope below $10^{-14}\,\si{\per\second}$ in $|L|/|L_{0}|$ is the
working acceptance threshold for a multi-decade mission. The
discipline that bounds the failure has three layers: a
symplectic or geometric integrator at the inner loop, a
runtime conservation monitor with an automatic trigger to
extended precision or step reduction, and a periodic
recalibration against an independent reference (star tracker,
GPS, deep-space-network range-rate). The reference layer
catches not only integrator drift but also actuator wear,
sensor drift, and modelling error in the dynamical model
itself.
\end{exercise}

\begin{exercise}
A control system for a chemical reactor is designed using a
linearised model around the steady-state operating point. After
deployment, the reactor experiences a transient that takes the
state into a regime where the linearisation is poor, and the
control system fails to restore the steady state. Argue, in
200 to 300 words, the failure mechanism in terms of the
phase-plane structure: the linearisation captured local
behaviour around one equilibrium; the transient pushed the
state past the basin of attraction. Identify the design
discipline that bounds the failure (controller saturation
analysis, basin-of-attraction estimation, fail-safe state
detection).

\paragraph{Discussion.} The linearised model captures the
Jacobian at the operating point: the local trace and
determinant, the local eigenvalues, the local time constants.
Nonlinear chemical reactors typically admit multiple
equilibria. A common pattern in exothermic reactors is two
stable equilibria (a low-conversion stable point and a
high-conversion stable point) separated by an unstable
intermediate (a saddle); the phase-plane picture has two
basins of attraction with the saddle's stable manifold as the
separatrix. The linear controller, designed for one
equilibrium, drives the state back along the local
linearisation; as long as the state remains within the basin,
it works. A transient that crosses the separatrix takes the
state into the other basin, where the linear control law no
longer points toward the design equilibrium; the controller
acts on the wrong attractor, and the reactor stabilises
elsewhere (or, if the second equilibrium is itself unstable
under the linear control law, runs away). The failure mechanism
is a basin-crossing event in a nonlinear phase plane that the
linear design did not see. The bounding discipline has three
layers. First, controller saturation analysis: every actuator
has a maximum rate and limit, and the analyst must verify that
the linear control law's actuation stays within those limits
across the design's operating envelope, including transients.
Second, basin-of-attraction estimation: Lyapunov-function
bounds or numerical sweeping characterise the design basin
explicitly, so the operating envelope can be sized below the
boundary with margin. Third, fail-safe state detection: a
runtime supervisor that monitors the state against the
characterised basin and triggers a safe-shutdown procedure on
crossing.
\end{exercise}

\subsection*{Judgment}\pathtag{standard}

\begin{exercise}
A junior engineer proposes to verify a pharmacokinetic ODE
simulation by running the integrator for $1{,}000\,\si{\hour}$
and observing whether the steady-state concentration matches the
analytical prediction. Write a one-paragraph response that takes
the proposal seriously, identifies the regime in which it is
informative (if the integrator is stable and accurate over that
window), and identifies the regime in which it is misleading (if
accumulated roundoff is dominant or the slow mode is sub-critically
sampled).

\paragraph{Discussion.} The proposal is sound as a smoke test
under three conditions. First, the integrator is stable for
the chosen step size, so the simulation does not blow up over
the $1{,}000\,\si{\hour}$ window. Second, the step size is fine
enough that the per-step truncation error is small compared to
the working tolerance (a one-percent tolerance on a $4\,\si{
\hour}$-half-life drug with $h = 0.1\,\si{\hour}$ gives
truncation $\sim 10^{-6}$ per step at fourth order). Third, the
accumulated roundoff over $N = 10^{4}$ steps stays below the
working tolerance: $\sqrt{N} \epsilon \sim 10^{2} \cdot 2.2
\times 10^{-16} \sim 10^{-14}$ at double precision, safely
negligible. Under those conditions the steady-state match is a
useful verification: it confirms both the analytical formula
and the integrator's long-term tracking. The proposal is
misleading in three regimes. First, if the integrator is
marginally stable (an explicit method on a problem near its
stability boundary), the simulation can drift over $1{,}000\,
\si{\hour}$ to a wrong steady state that happens to look
plausible. Second, if accumulated roundoff is dominant (an
incorrect step size, a low-precision arithmetic, or a
poorly-conditioned right-hand side), the steady-state value is
contaminated by arithmetic and the verification fails to
discriminate model from numerics. Third, if the slow mode is
sub-critically sampled (the dosing interval not resolved by the
chosen step, or a beat frequency between the dose cycle and the
step size), the integrator can produce systematic offset error
that mimics or masks the analytical prediction. The disciplined
response: run the proposal as a verification, but supplement it
with a step-size halving test, an extended-precision rerun, and
a sample-rate sensitivity sweep.
\end{exercise}

\begin{exercise}
Argue, in 200 to 300 words, when a closed-form analytical
solution of an ODE is preferable to a numerical solution, and
when the inverse holds. Take a position on a class of system
(linear time-invariant with constant forcing, weakly nonlinear
near a stable equilibrium, strongly nonlinear far from
equilibrium, stiff multi-scale) and identify which approach the
working engineer reaches for first.

\paragraph{Discussion.} The closed-form solution wins when the
system is linear time-invariant with constant or simple
sinusoidal forcing, or separable, or otherwise reducible to
quadrature. The closed form gives the parameter dependence
exactly: the engineer reads $\omega_{0} = \sqrt{k/m}$ and knows
how the natural frequency scales with stiffness and mass; a
numerical solution gives a value, not a scaling law. The closed
form supports sensitivity analysis, dimensional analysis, and
limiting-regime checks that numerical solutions support only
through repeated runs. For first-order linear ODEs, second-order
linear constant-coefficient ODEs with sinusoidal forcing, and
homogeneous linear systems, the closed form is the working
answer. The numerical solution wins when the system is nonlinear,
or has time-varying coefficients, or has a forcing function with
no analytical antiderivative, or is high-dimensional. The
strongly-nonlinear case far from equilibrium is the classical
numerical-solution regime: no closed form exists, perturbation
series diverge, and the answer is whatever the integrator
produces under disciplined error control. For weakly nonlinear
systems near a stable equilibrium, the working engineer reaches
for both: a closed-form linearisation for parameter dependence
and scaling, plus a numerical solution to check the linearisation
boundary. For stiff multi-scale systems, the implicit numerical
solution is the only practical answer; closed forms exist for
toy problems but are inaccessible at industrial scale. The
working position the engineer takes: the closed form is the
first reach when one exists; the numerical solution is the
default when it does not, supplemented by every closed-form
check available in the limiting regimes. The discipline is to
recognise when each tool applies; the failure mode is to use
one when the other is required.
\end{exercise}

\begin{exercise}
A first-year analyst proposes to model a population whose growth
rate clearly varies with season as a single-rate exponential
$\dot{N} = r N$ with $r$ chosen as the annual average. Write a
two-paragraph response identifying the regime in which the
average-rate model is informative and the regime in which the
seasonal forcing must be retained. Identify the diagnostic that
distinguishes the two cases.

\paragraph{Discussion.} The averaged-rate model is informative
when the question concerns multi-year trends and the seasonal
forcing is approximately periodic with mean equal to the
annual average. For a linear ODE $\dot{N} = r(t) N$ with
$r(t) = \bar{r} + r_{1}(t)$ and $r_{1}(t)$ a zero-mean
periodic function, the solution is $N(t) = N_{0} \exp(\bar{r} t
+ \int_{0}^{t} r_{1}(s) ds)$. The seasonal forcing produces a
bounded periodic modulation of the multi-year trend: the
modulation amplitude is set by $\max |\int r_{1}|$, the trend
by $\bar{r} t$. A questioner who cares about a five-year
population growth and is willing to accept a percent-level
seasonal modulation can use the averaged model; the working
trend is captured.

The averaged-rate model fails in three regimes. First, if the
seasonal $r_{1}(t)$ takes large positive and negative values,
the modulation amplitude is no longer small and the averaged
model misses the within-year peaks and troughs that matter for
management decisions. Second, if the system is nonlinear (a
logistic with seasonal $r$ or a saturating predation rate), the
linear averaging argument fails: the answer depends on the
forcing waveform, not just its mean, through the nonlinear
coupling between $r(t)$ and $N(t)$. Third, if a resonant
coupling exists (a predator-prey system whose intrinsic period
matches the seasonal forcing period), parametric resonance can
produce qualitative dynamics absent from the averaged model.
Diagnostic: compute the seasonal amplitude
$\sigma_{r} = \mathrm{std}(r(t))$ relative to the mean $\bar{r}$
and the ratio of seasonal to intrinsic period. If $\sigma_{r}/
|\bar{r}| < 0.2$ and the dynamics are linear with no resonance,
the averaged model is informative; otherwise, the seasonal
forcing must be retained.
\end{exercise}
